\documentclass{book}
\usepackage[utf8]{inputenc}
\usepackage[T1]{fontenc}
\usepackage{lmodern}
\usepackage[spanish]{babel}
\addto\captionsspanish{\renewcommand{\contentsname}{Índice}}
\AtBeginDocument{\spanishdeactivate{"~<>}}
\addto\extrasspanish{\spanishdeactivate{"~<>}}
\usepackage[a4paper,margin=2.5cm]{geometry}
\usepackage{titlesec}
\titleformat{\chapter}[display]
  {\normalfont\filright}{\large\scshape\chaptertitlename\ \thechapter}{1ex}{\huge\bfseries}
\titlespacing*{\chapter}{0pt}{2.5ex plus 1ex minus .2ex}{2.3ex}
\titleformat{name=\chapter,numberless}[display]
  {\normalfont\filright}{}{0pt}{\huge\bfseries}
\titleformat{\section}{\Large\bfseries}{{\thesection}}{0.8em}{}
\titleformat{\subsection}{\large\bfseries}{{\thesubsection}}{0.7em}{}
\setcounter{secnumdepth}{3}
\setcounter{tocdepth}{2}
\usepackage{graphicx}
\usepackage{tikz}
\usetikzlibrary{arrows.meta,positioning,calc,shapes.geometric}
\usepackage{xcolor}
\usepackage{amsmath}
\usepackage{amssymb}
\usepackage{array}
\usepackage{longtable}
\newcounter{studiumfigure}
\definecolor{studiumInk}{RGB}{28,40,58}
\definecolor{studiumRule}{RGB}{28,40,58}
\definecolor{studiumTipBack}{RGB}{232,244,236}
\definecolor{studiumTipFrame}{RGB}{27,107,58}
\definecolor{studiumDefBack}{RGB}{232,240,250}
\definecolor{studiumDefFrame}{RGB}{24,74,130}
\definecolor{studiumWorkedBack}{RGB}{255,244,230}
\definecolor{studiumWorkedFrame}{RGB}{138,78,12}
\definecolor{studiumCheckBack}{RGB}{243,236,248}
\definecolor{studiumCheckFrame}{RGB}{96,48,130}
\definecolor{studiumCodeBack}{RGB}{245,245,242}
\definecolor{studiumCodeRule}{RGB}{48,48,52}
\usepackage{listings}
\lstdefinelanguage{Rust}{
  morekeywords={as,async,await,break,const,continue,crate,dyn,else,enum,extern,false,fn,for,if,impl,in,let,loop,match,mod,move,mut,pub,ref,return,self,Self,static,struct,super,trait,true,type,unsafe,use,where,while},
  sensitive=true,
  morecomment=[l]{//},
  morecomment=[s]{/*}{*/},
  morestring=[b]{"}
}
\lstdefinelanguage{Go}{
  morekeywords={break,case,chan,const,continue,default,defer,else,fallthrough,for,func,go,goto,if,import,interface,map,package,range,return,select,struct,switch,type,var},
  sensitive=true,
  morecomment=[l]{//},
  morecomment=[s]{/*}{*/},
  morestring=[b]{"}
}
\lstdefinelanguage{JavaScript}{
  morekeywords={break,case,catch,class,const,continue,debugger,default,delete,do,else,export,extends,finally,for,function,if,import,in,instanceof,let,new,return,super,switch,this,throw,try,typeof,var,void,while,with,yield},
  sensitive=true,
  morecomment=[l]{//},
  morecomment=[s]{/*}{*/},
  morestring=[b]{"},
  morestring=[b]{'},
}
\lstdefinestyle{studium}{
  basicstyle=\ttfamily\small,
  columns=fullflexible,
  keepspaces=true,
  showstringspaces=false,
  breaklines=true,
  breakatwhitespace=false,
  frame=leftline,
  framerule=1.15pt,
  rulecolor=\color{studiumCodeRule},
  backgroundcolor=\color{studiumCodeBack},
  xleftmargin=1.15em,
  framexleftmargin=0.75em,
  aboveskip=0.9em,
  belowskip=0.7em,
  tabsize=4,
  keywordstyle=\ttfamily\bfseries,
  commentstyle=\ttfamily\itshape,
  stringstyle=\ttfamily
}
\lstset{style=studium}
\usepackage[breakable,skins]{tcolorbox}
\tcbset{
  studiumtip/.style={colback=studiumTipBack,colframe=studiumTipFrame,colbacktitle=studiumTipFrame,coltitle=white,coltext=black,fonttitle=\bfseries,boxrule=0.65pt,arc=0.7pt,left=1.7mm,right=1.7mm,top=1.3mm,bottom=1.3mm,before skip=10pt,after skip=10pt},
  studiumdef/.style={colback=studiumDefBack,colframe=studiumDefFrame,colbacktitle=studiumDefFrame,coltitle=white,coltext=black,fonttitle=\bfseries,boxrule=0.65pt,arc=0.7pt,left=1.7mm,right=1.7mm,top=1.3mm,bottom=1.3mm,before skip=10pt,after skip=10pt},
  studiumworked/.style={colback=studiumWorkedBack,colframe=studiumWorkedFrame,colbacktitle=studiumWorkedFrame,coltitle=white,coltext=black,fonttitle=\bfseries,boxrule=0.65pt,arc=0.7pt,left=1.7mm,right=1.7mm,top=1.3mm,bottom=1.3mm,before skip=10pt,after skip=10pt},
  studiumcheck/.style={colback=studiumCheckBack,colframe=studiumCheckFrame,colbacktitle=studiumCheckFrame,coltitle=white,coltext=black,fonttitle=\bfseries,boxrule=0.65pt,arc=0.7pt,left=1.7mm,right=1.7mm,top=1.3mm,bottom=1.3mm,before skip=10pt,after skip=10pt}
}
\titleformat{\section}{\Large\bfseries}{\thesection}{0.8em}{}
\titleformat{\subsection}{\large\bfseries}{\thesubsection}{0.7em}{}
\setcounter{secnumdepth}{3}
\setcounter{tocdepth}{2}
\usepackage{fancyhdr}
\pagestyle{fancy}
\fancyhf{}
\fancyhead[LE]{\small\thepage\hspace{1.5em}\scshape\nouppercase{\leftmark}}
\fancyhead[RO]{\small\scshape\nouppercase{\rightmark}\hspace{1.5em}\thepage}
\fancyfoot[C]{\small Borrador}
\renewcommand{\headrulewidth}{0.4pt}
\renewcommand{\footrulewidth}{0pt}
\fancypagestyle{plain}{%
  \fancyhf{}
  \fancyfoot[C]{\small Borrador}
  \renewcommand{\headrulewidth}{0pt}
  \renewcommand{\footrulewidth}{0pt}
}
\usepackage[hyphens]{url}
\usepackage[hidelinks]{hyperref}
\setlength{\emergencystretch}{1.5em}
\makeatletter
\@openrightfalse
\let\cleardoublepage\clearpage
\makeatother
\begin{document}
\frontmatter
\thispagestyle{empty}
\setlength{\parindent}{0pt}
\vspace*{0.08\textheight}
{\normalsize\scshape\color{studiumInk} Studium\par}
{\small Enciclopedia de Programación e Ingeniería del Software\par}
\vspace{0.55em}
{\color{studiumRule}\rule{\linewidth}{0.9pt}\par}
\vspace{0.16\textheight}
{\raggedright\Huge\bfseries Volumen VI: Monografías de C\#, Java, Go, Kotlin y Swift\par}
\vspace{0.9em}
{\raggedright\Large\color{studiumInk} Parte V (Bloques A\(-\)H): Máquinas virtuales gestionadas (CLR y JVM), concurrencia G-M-P en Go, Coroutines y ARC/Actores\par}
\vfill
{\color{studiumRule}\rule{0.36\linewidth}{0.45pt}\par}
\vspace{1.05em}
{\large\scshape Borrador\par}
\vspace{0.45em}
{\normalsize Edición 0.1.0 --- No publicado\par}
\vspace{0.4em}
{\normalsize\today\par}
\vspace*{0.08\textheight}
\clearpage
\chapter{Prefacio}
\noindent Este libro presenta el tema con el lenguaje del curso. Cada capítulo abre con una entrada breve, sigue con la explicación y cierra con un problema resuelto y una autoficha.
\chapter{Cómo usar este libro}
\noindent Lee el capítulo seguido. El consejo, la definición, el problema resuelto y la autoficha van en recuadros. La explicación es el cuerpo del texto.

\noindent Un capítulo lleva una entrada en cursiva, la explicación en el cuerpo, como mucho un consejo, definiciones solo al introducir un término, un problema resuelto y una autoficha.
\tableofcontents
\mainmatter

\part{Monografías de Plataformas Gestionadas: C\# (.NET 8/9) y Java (JVM 21 LTS)}

\chapter{Monografía de C\# (I \(-\) Bloques A\(-\)D): Arquitectura .NET/CLR, Tipos Valor vs Referencia, Records, Pattern Matching y LINQ}
\label{chap:cap01-csharp-dotnet-clr-tipos-valor-referencia-linq}
% studium:record eyJpZCI6ICJQQVItMDAwMSIsICJzZWN0aW9uIjogImNhcDAxLWNzaGFycC1kb3RuZXQtY2xyLXRpcG9zLXZhbG9yLXJlZmVyZW5jaWEtbGlucSJ9

\noindent\textit{Este capítulo inicia la monografía autónoma de **C\# (.NET 8/9)** (Bloques A, B, C y D). Diseñado por Anders Hejlsberg, C\# combina la ergonomía expresiva de los lenguajes funcionales modernos con un control de bajo nivel sobre la disposición de memoria (\texttt{struct}, \texttt{Span<T>}, genéricos reificados en la máquina virtual) único entre los entornos gestionados.}
% studium:end-record eyJpZCI6ICJQQVItMDAwMSIsICJzZWN0aW9uIjogImNhcDAxLWNzaGFycC1kb3RuZXQtY2xyLXRpcG9zLXZhbG9yLXJlZmVyZW5jaWEtbGlucSJ9

\begin{tcolorbox}[title={Consejo}, studiumtip, breakable]
% studium:record eyJpZCI6ICJQQVItMDAwNSIsICJzZWN0aW9uIjogImNhcDAxLWNzaGFycC1kb3RuZXQtY2xyLXRpcG9zLXZhbG9yLXJlZmVyZW5jaWEtbGlucSJ9

Cuando declares un \texttt{struct} en C\#, márcalo como \texttt{readonly struct} (o \texttt{readonly record struct}) siempre que sea posible: si pasas un \texttt{struct} no \texttt{readonly} mediante \texttt{in} o lo almacenas en un campo \texttt{readonly}, el compilador se ve obligado a emitir copias defensivas ocultas en cada llamada a método para proteger el original.
% studium:end-record eyJpZCI6ICJQQVItMDAwNSIsICJzZWN0aW9uIjogImNhcDAxLWNzaGFycC1kb3RuZXQtY2xyLXRpcG9zLXZhbG9yLXJlZmVyZW5jaWEtbGlucSJ9
\end{tcolorbox}

\begin{tcolorbox}[title={Definición}, studiumdef, breakable]
% studium:record eyJpZCI6ICJQQVItMDAwNCIsICJzZWN0aW9uIjogImNhcDAxLWNzaGFycC1kb3RuZXQtY2xyLXRpcG9zLXZhbG9yLXJlZmVyZW5jaWEtbGlucSJ9

Reificación de Genéricos en el CLR (.NET): preservación de los argumentos de tipo genérico en los metadatos CIL y en tiempo de ejecución, permitiendo que el JIT emita representaciones especializadas en memoria para cada tipo por valor (\texttt{struct}) sin encajonamiento (*boxing*).
% studium:end-record eyJpZCI6ICJQQVItMDAwNCIsICJzZWN0aW9uIjogImNhcDAxLWNzaGFycC1kb3RuZXQtY2xyLXRpcG9zLXZhbG9yLXJlZmVyZW5jaWEtbGlucSJ9
\end{tcolorbox}

\begin{tcolorbox}[title={Autoficha}, studiumcheck, breakable]
% studium:record eyJpZCI6ICJQQVItMDAwNiIsICJzZWN0aW9uIjogImNhcDAxLWNzaGFycC1kb3RuZXQtY2xyLXRpcG9zLXZhbG9yLXJlZmVyZW5jaWEtbGlucSJ9

¿Qué problema grave de rendimiento ocurre en una base de datos de millones de filas si un desarrollador convierte un \texttt{IQueryable<Order>} de EF Core a \texttt{IEnumerable<Order>} mediante \texttt{.AsEnumerable()} o \texttt{.ToList()} **antes** de aplicar \texttt{.Where(...)}?
% studium:end-record eyJpZCI6ICJQQVItMDAwNiIsICJzZWN0aW9uIjogImNhcDAxLWNzaGFycC1kb3RuZXQtY2xyLXRpcG9zLXZhbG9yLXJlZmVyZW5jaWEtbGlucSJ9
\end{tcolorbox}

\section{Bloques A\(-\)C: Historia de C\# y .NET, CLI (`dotnet`), CIL, tipos valor vs referencia y nulabilidad}
% studium:record eyJpZCI6ICJQQVItMDAwMiIsICJzZWN0aW9uIjogImNhcDAxLWNzaGFycC1kb3RuZXQtY2xyLXRpcG9zLXZhbG9yLXJlZmVyZW5jaWEtbGlucSJ9

En la plataforma **.NET**, el compilador Roslyn (\texttt{csc}) traduce el código fuente C\# a *Common Intermediate Language* (CIL, estándar ECMA-335) almacenado en ensamblados \texttt{.dll}. Durante la ejecución, el compilador JIT **RyuJIT** traduce los métodos CIL a código máquina nativo aplicando compilación por niveles (*Tiered Compilation*) y perfilado dinámico (Dynamic PGO), o bien compila directamente a binario nativo autónomo mediante **Native AOT**. Una decisión crucial de ingeniería en C\# es la distinción entre **tipos por referencia** (\texttt{class}, alojados en el montículo gestionado con una cabecera de objeto de 16 bytes en 64 bits) y **tipos por valor** (\texttt{struct}, alojados en línea en el marco de pila o directamente dentro del array/objeto contenedor sin indirección de puntero). Además, a diferencia de la JVM, el CLR **reifica los tipos genéricos**: \texttt{List<int>} o \texttt{Dictionary<long, double>} generan código máquina especializado con arrays planos de valores primitivos sin *boxing*.
% studium:end-record eyJpZCI6ICJQQVItMDAwMiIsICJzZWN0aW9uIjogImNhcDAxLWNzaGFycC1kb3RuZXQtY2xyLXRpcG9zLXZhbG9yLXJlZmVyZW5jaWEtbGlucSJ9

\section{Bloque D: Propiedades, records, interfaces, genéricos reificados, delegados y LINQ}
% studium:record eyJpZCI6ICJQQVItMDAwMyIsICJzZWN0aW9uIjogImNhcDAxLWNzaGFycC1kb3RuZXQtY2xyLXRpcG9zLXZhbG9yLXJlZmVyZW5jaWEtbGlucSJ9

C\# moderno incorpora tipos de referencia anulables verificados por flujo (\texttt{<Nullable>enable</Nullable>}), tipos de datos inmutables por valor o referencia (\texttt{readonly record struct} y \texttt{record class}) con igualdad estructural sintetizada y **LINQ** (*Language Integrated Query*). En LINQ es indispensable distinguir entre \texttt{IEnumerable<T>}, cuyos operadores reciben delegados \texttt{Func<T, bool>} evaluados en memoria en el proceso actual, e \texttt{IQueryable<T>}, cuyos operadores reciben árboles de expresión (\texttt{Expression<Func<T, bool>>}) que un proveedor como Entity Framework Core inspecciona para traducirlos a consultas SQL ejecutadas en el motor de base de datos:
\begin{lstlisting}[style=studium]
public readonly record struct SensorSample(int SensorId, double Value, long TimestampMs);

public static double ComputeHighValueAverage(ReadOnlySpan<SensorSample> samples, double threshold)
{
    double sum = 0.0;
    int count = 0;
    foreach (ref readonly var s in samples)
    {
        if (s.Value >= threshold) { sum += s.Value; count++; }
    }
    return count > 0 ? sum / count : 0.0;
}
\end{lstlisting}
Observa cómo \texttt{foreach (ref readonly var s in samples)} sobre un \texttt{ReadOnlySpan<SensorSample>} recorre las estructuras contiguas por referencia constante sin realizar ni una sola copia ni alojar un iterador en el montículo.
% studium:end-record eyJpZCI6ICJQQVItMDAwMyIsICJzZWN0aW9uIjogImNhcDAxLWNzaGFycC1kb3RuZXQtY2xyLXRpcG9zLXZhbG9yLXJlZmVyZW5jaWEtbGlucSJ9

\begin{tcolorbox}[title={Problema resuelto}, studiumworked, breakable]
% studium:record eyJpZCI6ICJQUkItMDAwMSIsICJzZWN0aW9uIjogImNhcDAxLWNzaGFycC1kb3RuZXQtY2xyLXRpcG9zLXZhbG9yLXJlZmVyZW5jaWEtbGlucSJ9

\noindent\textbf{Enunciado}
En una aplicación financiera en C\# sobre .NET de 64 bits, se mantiene un búfer de \(N = 100000\) cotizaciones donde cada cotización tiene \(S_{\text{val}} = 16\text{ B}\) de datos (\texttt{long Timestamp}, \texttt{double Price}). Si se modela como \texttt{readonly record struct}, el array \texttt{Quote[]} ocupa \(100000 \times 16\text{ B}\). Si se modelara como \texttt{class}, cada elemento añadiría \(8\text{ B}\) de puntero en el array más \(16\text{ B}\) de cabecera de objeto en el montículo (\(24\text{ B}\) extra por instancia). Calcula cuántos kilobytes (dividiendo entre 1000) ahorra el diseño con \texttt{struct}.

\noindent\textbf{Datos}
\begin{itemize}
\item $N = 100000$ — Número de cotizaciones en el array
\item $S_{extra} = 24\,\mathrm{B}$ — Sobrecoste por instancia de class (8 B ref + 16 B cabecera)
\end{itemize}

\noindent\textbf{Incógnita}
Magnitud objetivo solicitada en el enunciado

\noindent\textbf{Modelo físico}

\noindent\textit{Ahorro de memoria entre array de struct y array de class en el CLR (.NET)}

\[
\Delta M_{\text{kB}} = \frac{100000 \cdot 24}{1000}
\]

\noindent\textbf{Hipótesis}
\begin{itemize}
\item Evaluación determinista bajo el modelo formal del capítulo
\end{itemize}

\noindent\textbf{Desarrollo}
\begin{enumerate}
\item
Multiplicamos las 100000 instancias por 24 B y dividimos entre 1000 B/kB

\[
\Delta M_{\text{kB}} = \frac{100000 \cdot 24}{1000}
\]

\[
\frac{\left(100000\cdot 24\right)}{1000}
\]
\end{enumerate}

\noindent\textbf{Resultado}

\[
Resultado = 2400
\]

\noindent\textbf{Interpretación}
El diseño con readonly record struct ahorra 2400 kB (2.4 MB) y elimina 100 000 nodos del grafo del recolector de basura.

\noindent\textbf{Comprobación de los pasos}
\begin{itemize}
\item Paso 1: $2400$ — reproducido
\end{itemize}

\noindent\textbf{Comprobación de la respuesta} la aritmética del último paso reproduce la respuesta enunciada a la precisión con que está escrita. Es una reproducción aritmética; no valida el modelo físico ni sus hipótesis.
% studium:end-record eyJpZCI6ICJQUkItMDAwMSIsICJzZWN0aW9uIjogImNhcDAxLWNzaGFycC1kb3RuZXQtY2xyLXRpcG9zLXZhbG9yLXJlZmVyZW5jaWEtbGlucSJ9
\end{tcolorbox}

\begin{tcolorbox}[title={Ejercicios}, studiumworked, breakable]

\noindent\textbf{1 [PRB-0003] (Básico, Numérico)}
En .NET de 64 bits, la cabecera de todo objeto por referencia (\texttt{class}) consta de \texttt{ObjHeader} (8 B) y \texttt{MethodTable*} (8 B). ¿Cuántos bytes suma esta cabecera?

\noindent\textbf{2 [PRB-0004] (Básico, Numérico)}
Un \texttt{readonly record struct Point3D(double X, double Y, double Z)} contiene tres \texttt{double} de 8 bytes. ¿Cuántos bytes ocupa cada instancia?

\noindent\textbf{3 [PRB-0005] (Intermedio, Desarrollo por pasos)}
Una consulta LINQ encadena \texttt{.Where(...)} y \texttt{.Take(25)} sobre una secuencia perezosa de 10 000 elementos donde se cumple el predicado en 1 de cada 4 elementos. ¿Cuántos elementos evalúa antes de detenerse (\(25 \times 4\))?

\noindent\textbf{4 [PRB-0006] (Intermedio, Diseño)}
El recolector generacional del CLR divide el montículo de objetos pequeños en 3 generaciones (\texttt{Gen 0}, \texttt{Gen 1}, \texttt{Gen 2}). Si \texttt{Gen 0} tiene un presupuesto de 4 MiB y \texttt{Gen 1} de 2 MiB, ¿cuántos KiB suman las dos generaciones efímeras (\((4 + 2) \times 1024\))?
\end{tcolorbox}

\chapter{Monografía de C\# (II \(-\) Bloques E\(-\)H): Span<T>, Memory<T>, Async/Await, ValueTask, ASP.NET Core y Proyectos}
\label{chap:cap02-csharp-avanzado-experto-span-async-valuetask-aspnet-proyectos}
% studium:record eyJpZCI6ICJQQVItMDAwNyIsICJzZWN0aW9uIjogImNhcDAyLWNzaGFycC1hdmFuemFkby1leHBlcnRvLXNwYW4tYXN5bmMtdmFsdWV0YXNrLWFzcG5ldC1wcm95ZWN0b3MifQ==

\noindent\textit{Este capítulo completa la monografía de C\# cubriendo los niveles Avanzado y Experto (Bloques E y F), el Ecosistema empresarial (Bloque G) y los Proyectos (Bloque H). Explicamos cómo las primitivas \texttt{Span<T>} y \texttt{ValueTask<T>} han situado a ASP.NET Core (Kestrel) entre los servidores web más rápidos del mundo.}
% studium:end-record eyJpZCI6ICJQQVItMDAwNyIsICJzZWN0aW9uIjogImNhcDAyLWNzaGFycC1hdmFuemFkby1leHBlcnRvLXNwYW4tYXN5bmMtdmFsdWV0YXNrLWFzcG5ldC1wcm95ZWN0b3MifQ==

\begin{tcolorbox}[title={Consejo}, studiumtip, breakable]
% studium:record eyJpZCI6ICJQQVItMDAxMSIsICJzZWN0aW9uIjogImNhcDAyLWNzaGFycC1hdmFuemFkby1leHBlcnRvLXNwYW4tYXN5bmMtdmFsdWV0YXNrLWFzcG5ldC1wcm95ZWN0b3MifQ==

En todos los métodos asíncronos de entrada/salida en C\#, acepta siempre un \texttt{CancellationToken ct = default} como último parámetro y pásalo hacia abajo a todas las llamadas de red y base de datos para abortar el trabajo inmediatamente cuando el cliente HTTP cierra la conexión.
% studium:end-record eyJpZCI6ICJQQVItMDAxMSIsICJzZWN0aW9uIjogImNhcDAyLWNzaGFycC1hdmFuemFkby1leHBlcnRvLXNwYW4tYXN5bmMtdmFsdWV0YXNrLWFzcG5ldC1wcm95ZWN0b3MifQ==
\end{tcolorbox}

\begin{tcolorbox}[title={Definición}, studiumdef, breakable]
% studium:record eyJpZCI6ICJQQVItMDAxMCIsICJzZWN0aW9uIjogImNhcDAyLWNzaGFycC1hdmFuemFkby1leHBlcnRvLXNwYW4tYXN5bmMtdmFsdWV0YXNrLWFzcG5ldC1wcm95ZWN0b3MifQ==

Estructura restringida a pila (\texttt{ref struct} en C\#): tipo por valor cuya ubicación está garantizada estáticamente por el compilador y el verificador del CLR exclusivamente en la pila de ejecución, permitiendo contener punteros interiores gestionados (\texttt{ref T} en \texttt{Span<T>}) seguros frente a compactaciones del recolector de basura.
% studium:end-record eyJpZCI6ICJQQVItMDAxMCIsICJzZWN0aW9uIjogImNhcDAyLWNzaGFycC1hdmFuemFkby1leHBlcnRvLXNwYW4tYXN5bmMtdmFsdWV0YXNrLWFzcG5ldC1wcm95ZWN0b3MifQ==
\end{tcolorbox}

\begin{tcolorbox}[title={Autoficha}, studiumcheck, breakable]
% studium:record eyJpZCI6ICJQQVItMDAxMiIsICJzZWN0aW9uIjogImNhcDAyLWNzaGFycC1hdmFuemFkby1leHBlcnRvLXNwYW4tYXN5bmMtdmFsdWV0YXNrLWFzcG5ldC1wcm95ZWN0b3MifQ==

¿Por qué un objeto \texttt{Span<byte>} no puede ser un campo de una clase ni persistir a través de una expresión \texttt{await} en un método \texttt{async}?
% studium:end-record eyJpZCI6ICJQQVItMDAxMiIsICJzZWN0aW9uIjogImNhcDAyLWNzaGFycC1hdmFuemFkby1leHBlcnRvLXNwYW4tYXN5bmMtdmFsdWV0YXNrLWFzcG5ldC1wcm95ZWN0b3MifQ==
\end{tcolorbox}

\section{Bloques E y F: Span<T>, ref struct, ArrayPool<T>, máquinas de estados async/await, ValueTask<T> y Channels}
% studium:record eyJpZCI6ICJQQVItMDAwOCIsICJzZWN0aW9uIjogImNhcDAyLWNzaGFycC1hdmFuemFkby1leHBlcnRvLXNwYW4tYXN5bmMtdmFsdWV0YXNrLWFzcG5ldC1wcm95ZWN0b3MifQ==

Para permitir el acceso uniforme y sin copias a memoria gestionada (\texttt{T[]}), memoria en pila (\texttt{stackalloc}) y memoria nativa no gestionada, C\# 7.2+ introdujo \texttt{Span<T>} y \texttt{ReadOnlySpan<T>}. Un \texttt{Span<T>} se declara como \texttt{ref struct}, lo que significa que el compilador y el verificador del CLR prohíben estáticamente que escape al montículo (no puede ser un campo de una \texttt{class}, ni ser encajonado como \texttt{object}, ni capturarse en una lambda, ni cruzar un punto \texttt{await}). En programación asíncrona, cada método \texttt{async} es transformado por el compilador en una estructura que implementa \texttt{IAsyncStateMachine}. Si un método \texttt{async Task<T>} completa su trabajo de forma asíncrona, aloja un objeto \texttt{Task<T>} en el montículo; pero en rutas calientes (como lecturas desde caché o sockets con datos ya presentes en el búfer del kernel) donde el 90 \% de las llamadas se resuelven síncronamente, devolver \texttt{ValueTask<T>} (un \texttt{readonly struct}) evita por completo cualquier asignación en el montículo:
\begin{lstlisting}[style=studium,escapeinside={(*@}{@*)}]
using System.Buffers;
using System.Threading;
using System.Threading.Tasks;

public sealed class CachedMetricReader
{
    private double? _cached;
    public ValueTask<double> ReadMetricAsync(CancellationToken ct)
    {
        if (_cached is double ready)
            return new ValueTask<double>(ready); // 0 bytes alojados en el mont(*@{\'i}@*)culo
        return new ValueTask<double>(FetchSlowAsync(ct));
    }
    private async Task<double> FetchSlowAsync(CancellationToken ct)
    {
        await Task.Delay(10, ct).ConfigureAwait(false);
        _cached = 42.5;
        return 42.5;
    }
}
\end{lstlisting}
Recuerda la restricción contractual de \texttt{ValueTask<T>}: nunca debe ser esperado (\texttt{await}) más de una vez ni consumido concurrentemente antes de que complete.
% studium:end-record eyJpZCI6ICJQQVItMDAwOCIsICJzZWN0aW9uIjogImNhcDAyLWNzaGFycC1hdmFuemFkby1leHBlcnRvLXNwYW4tYXN5bmMtdmFsdWV0YXNrLWFzcG5ldC1wcm95ZWN0b3MifQ==

\section{Bloques G y H: Ecosistema (ASP.NET Core, DI, EF Core, xUnit, BenchmarkDotNet) y los 4 proyectos en C\#}
% studium:record eyJpZCI6ICJQQVItMDAwOSIsICJzZWN0aW9uIjogImNhcDAyLWNzaGFycC1hdmFuemFkby1leHBlcnRvLXNwYW4tYXN5bmMtdmFsdWV0YXNrLWFzcG5ldC1wcm95ZWN0b3MifQ==

En el Bloque G (Ecosistema), las aplicaciones .NET modernas se construyen sobre el contenedor de **Inyección de Dependencias (DI)** integrado en \texttt{Microsoft.Extensions.DependencyInjection}, que distingue tres ciclos de vida: \texttt{Transient} (una instancia nueva en cada resolución), \texttt{Scoped} (una instancia única por ámbito o petición HTTP, obligatoria para \texttt{DbContext} de Entity Framework Core) y \texttt{Singleton} (una única instancia segura para hilos durante toda la vida del proceso). Un antipatrón crítico (detectado automáticamente por \texttt{ValidateScopes = true} en entorno de desarrollo) es la **dependencia cautiva** (*Captive Dependency*): inyectar un servicio \texttt{Scoped} dentro de un \texttt{Singleton}, lo que retiene el servicio \texttt{Scoped} indefinidamente y provoca accesos concurrentes ilegales a \texttt{DbContext}.
% studium:end-record eyJpZCI6ICJQQVItMDAwOSIsICJzZWN0aW9uIjogImNhcDAyLWNzaGFycC1hdmFuemFkby1leHBlcnRvLXNwYW4tYXN5bmMtdmFsdWV0YXNrLWFzcG5ldC1wcm95ZWN0b3MifQ==

\begin{tcolorbox}[title={Problema resuelto}, studiumworked, breakable]
% studium:record eyJpZCI6ICJQUkItMDAwNyIsICJzZWN0aW9uIjogImNhcDAyLWNzaGFycC1hdmFuemFkby1leHBlcnRvLXNwYW4tYXN5bmMtdmFsdWV0YXNrLWFzcG5ldC1wcm95ZWN0b3MifQ==

\noindent\textbf{Enunciado}
Un servicio de alta frecuencia en ASP.NET Core procesa \(R = 50000\) llamadas por segundo donde el \(90\%\) (\(45000\) llamadas/s) se resuelven síncronamente desde memoria local y solo el \(10\%\) (\(5000\) llamadas/s) requieren E/S asíncrona real. Si cada objeto \texttt{Task<Result>} alojado en el montículo mide \(S_{\text{task}} = 72\text{ B}\), calcula cuántos kilobytes por segundo (dividiendo entre 1000) deja de alojar en \texttt{Gen 0} el servicio al cambiar el tipo de retorno de \texttt{Task<Result>} a \texttt{ValueTask<Result>}.

\noindent\textbf{Datos}
\begin{itemize}
\item $R_{sync} = 45000\,\mathrm{llamadas}\,\mathrm{s}^{-1}$ — Llamadas resueltas síncronamente por segundo
\item $S_{task} = 72\,\mathrm{B}$ — Tamaño del objeto Task<Result> evitado
\end{itemize}

\noindent\textbf{Incógnita}
Magnitud objetivo solicitada en el enunciado

\noindent\textbf{Modelo físico}

\noindent\textit{Ahorro de asignaciones en Gen 0 con ValueTask<T> en rutas calientes}

\[
\Delta A_{\text{kB/s}} = \frac{45000 \cdot 72}{1000}
\]

\noindent\textbf{Hipótesis}
\begin{itemize}
\item Evaluación determinista bajo el modelo formal del capítulo
\end{itemize}

\noindent\textbf{Desarrollo}
\begin{enumerate}
\item
Multiplicamos las 45000 llamadas síncronas por 72 B y dividimos entre 1000 B/kB

\[
\Delta A_{\text{kB/s}} = \frac{45000 \cdot 72}{1000}
\]

\[
\frac{\left(45000\cdot 72\right)}{1000}
\]
\end{enumerate}

\noindent\textbf{Resultado}

\[
Resultado = 3240
\]

\noindent\textbf{Interpretación}
Se evita alojar y recolectar 3240 kB/s (3.24 MB/s) de objetos efímeros en Gen 0.

\noindent\textbf{Comprobación de los pasos}
\begin{itemize}
\item Paso 1: $3240$ — reproducido
\end{itemize}

\noindent\textbf{Comprobación de la respuesta} la aritmética del último paso reproduce la respuesta enunciada a la precisión con que está escrita. Es una reproducción aritmética; no valida el modelo físico ni sus hipótesis.
% studium:end-record eyJpZCI6ICJQUkItMDAwNyIsICJzZWN0aW9uIjogImNhcDAyLWNzaGFycC1hdmFuemFkby1leHBlcnRvLXNwYW4tYXN5bmMtdmFsdWV0YXNrLWFzcG5ldC1wcm95ZWN0b3MifQ==
\end{tcolorbox}

\begin{tcolorbox}[title={Ejercicios}, studiumworked, breakable]

\noindent\textbf{1 [PRB-0009] (Básico, Numérico)}
En .NET de 64 bits, un \texttt{Span<byte>} contiene una referencia interior \texttt{ref byte} (8 B) y una longitud \texttt{int} alineada a 8 B. ¿Cuántos bytes ocupa en la pila?

\noindent\textbf{2 [PRB-0010] (Básico, Numérico)}
Si se alquila un búfer con \texttt{ArrayPool<byte>.Shared.Rent(4000)}, el pool devuelve un array cuya longitud es la siguiente potencia de dos (\(4096\)). ¿Cuántos bytes adicionales respecto a los 4000 solicitados tiene el array devuelto?

\noindent\textbf{3 [PRB-0011] (Intermedio, Desarrollo por pasos)}
Un canal acotado \texttt{Channel.CreateBounded<Job>(new BoundedChannelOptions(128))} contiene actualmente 95 trabajos. ¿Cuántos trabajos más pueden escribirse sin suspender al productor?

\noindent\textbf{4 [PRB-0012] (Intermedio, Diseño)}
Un contenedor de inyección de dependencias resuelve por cada petición HTTP 1 servicio \texttt{Singleton} (ya creado), 2 servicios \texttt{Scoped} y 4 servicios \texttt{Transient} (invocados 2 veces cada uno en la petición). ¿Cuántas instancias nuevas se crean durante esa petición (\(2 + 4 \times 2\))?
\end{tcolorbox}

\chapter{Monografía de Java (I \(-\) Bloques A\(-\)D): Arquitectura JVM, Bytecode, Memoria, Records, Sealed Classes, Generics y Streams}
\label{chap:cap03-java-jvm-bytecode-memoria-generics-streams}
% studium:record eyJpZCI6ICJQQVItMDAxMyIsICJzZWN0aW9uIjogImNhcDAzLWphdmEtanZtLWJ5dGVjb2RlLW1lbW9yaWEtZ2VuZXJpY3Mtc3RyZWFtcyJ9

\noindent\textit{Este capítulo inicia la monografía autónoma de **Java (Java 21 LTS)** (Bloques A, B, C y D). Java impulsa la infraestructura crítica financiera, bancaria y de comercio electrónico a escala planetaria, y con la llegada de Java 17 y Java 21 LTS ha incorporado tipos algebraicos de datos (\texttt{record} + \texttt{sealed interface}) e hilos virtuales que transforman por completo la escritura de código en la JVM.}
% studium:end-record eyJpZCI6ICJQQVItMDAxMyIsICJzZWN0aW9uIjogImNhcDAzLWphdmEtanZtLWJ5dGVjb2RlLW1lbW9yaWEtZ2VuZXJpY3Mtc3RyZWFtcyJ9

\begin{tcolorbox}[title={Consejo}, studiumtip, breakable]
% studium:record eyJpZCI6ICJQQVItMDAxNyIsICJzZWN0aW9uIjogImNhcDAzLWphdmEtanZtLWJ5dGVjb2RlLW1lbW9yaWEtZ2VuZXJpY3Mtc3RyZWFtcyJ9

Cuando operes con flujos numéricos en la API de Streams de Java, utiliza siempre las especializaciones primitivas (\texttt{IntStream}, \texttt{LongStream}, \texttt{DoubleStream} mediante \texttt{.mapToLong(...)} o \texttt{.mapToDouble(...)}) en lugar de \texttt{Stream<Integer>} o \texttt{Stream<Double>} para evitar el encajonamiento (*boxing*) de millones de objetos.
% studium:end-record eyJpZCI6ICJQQVItMDAxNyIsICJzZWN0aW9uIjogImNhcDAzLWphdmEtanZtLWJ5dGVjb2RlLW1lbW9yaWEtZ2VuZXJpY3Mtc3RyZWFtcyJ9
\end{tcolorbox}

\begin{tcolorbox}[title={Definición}, studiumdef, breakable]
% studium:record eyJpZCI6ICJQQVItMDAxNiIsICJzZWN0aW9uIjogImNhcDAzLWphdmEtanZtLWJ5dGVjb2RlLW1lbW9yaWEtZ2VuZXJpY3Mtc3RyZWFtcyJ9

Principio PECS (Producer Extends, Consumer Super en Java): regla de varianza en el sitio de uso (*use-site variance*) según la cual una estructura parametrizada que actúa como productora de valores \texttt{T} debe tiparse de forma covariante (\texttt{? extends T}), y una que actúa como consumidora de valores \texttt{T} de forma contravariante (\texttt{? super T}).
% studium:end-record eyJpZCI6ICJQQVItMDAxNiIsICJzZWN0aW9uIjogImNhcDAzLWphdmEtanZtLWJ5dGVjb2RlLW1lbW9yaWEtZ2VuZXJpY3Mtc3RyZWFtcyJ9
\end{tcolorbox}

\begin{tcolorbox}[title={Autoficha}, studiumcheck, breakable]
% studium:record eyJpZCI6ICJQQVItMDAxOCIsICJzZWN0aW9uIjogImNhcDAzLWphdmEtanZtLWJ5dGVjb2RlLW1lbW9yaWEtZ2VuZXJpY3Mtc3RyZWFtcyJ9

¿Por qué en Java \texttt{Integer.valueOf(100) == Integer.valueOf(100)} evalúa a \texttt{true} (por la caché obligatoria de \texttt{-128} a \texttt{127} de la JLS §5.1.7), pero \texttt{Integer.valueOf(1000) == Integer.valueOf(1000)} evalúa a \texttt{false} y debe compararse siempre con \texttt{.equals()}?
% studium:end-record eyJpZCI6ICJQQVItMDAxOCIsICJzZWN0aW9uIjogImNhcDAzLWphdmEtanZtLWJ5dGVjb2RlLW1lbW9yaWEtZ2VuZXJpY3Mtc3RyZWFtcyJ9
\end{tcolorbox}

\section{Bloques A\(-\)C: Evolución hasta Java 21 LTS, arquitectura JVM, JIT C1/C2, G1/ZGC y sintaxis fundamental}
% studium:record eyJpZCI6ICJQQVItMDAxNCIsICJzZWN0aW9uIjogImNhcDAzLWphdmEtanZtLWJ5dGVjb2RlLW1lbW9yaWEtZ2VuZXJpY3Mtc3RyZWFtcyJ9

Creado por James Gosling en Sun Microsystems en 1995, Java compila los ficheros \texttt{.java} mediante \texttt{javac} a código de bytes de pila (\texttt{.class}) independiente de la arquitectura física. En tiempo de ejecución, la máquina virtual **HotSpot** aplica compilación escalonada (*Tiered Compilation*): comienza interpretando el bytecode (Nivel 0), compila métodos frecuentes con el compilador rápido C1 recolectando perfiles de ejecución (Niveles \(1-3)\) y finalmente optimiza el código caliente con el compilador servidor **C2** (Nivel 4), capaz de desvirtualizar llamadas polimórficas monomórficas/bimórficas y aplicar análisis de escape (*Escape Analysis*) para alojar objetos efímeros en registros o escalarizarlos en la pila. Para la gestión del montículo, Java 21 ofrece **G1GC** (recolector regional por defecto) y **Generational ZGC** (\texttt{-XX:+UseZGC -XX:+ZGenerational}), un recolector concurrente con barreras de punteros coloreados (*colored pointers*) cuyas pausas *Stop-The-World* son inferiores a \(1\text{ ms}\) incluso en montículos de varios terabytes.
% studium:end-record eyJpZCI6ICJQQVItMDAxNCIsICJzZWN0aW9uIjogImNhcDAzLWphdmEtanZtLWJ5dGVjb2RlLW1lbW9yaWEtZ2VuZXJpY3Mtc3RyZWFtcyJ9

\section{Bloque D: Records, Sealed Classes, Pattern Matching, Genéricos (PECS), Colecciones y Streams API}
% studium:record eyJpZCI6ICJQQVItMDAxNSIsICJzZWN0aW9uIjogImNhcDAzLWphdmEtanZtLWJ5dGVjb2RlLW1lbW9yaWEtZ2VuZXJpY3Mtc3RyZWFtcyJ9

Java 17 y Java 21 LTS introdujeron el paradigma de **Programación Orientada a Datos** (*Data-Oriented Programming*): la combinación de clases \texttt{record} (portadores transparentes e inmutables de datos), interfaces selladas (\texttt{sealed interface ... permits ...}) y coincidencia de patrones exhaustiva en expresiones \texttt{switch} con desestructuración:
\begin{lstlisting}[style=studium,language={Java}]
public sealed interface PaymentResult permits PaymentResult.Approved, PaymentResult.Declined {
    record Approved(String txId, long cents) implements PaymentResult {}
    record Declined(String reason, int retryAfterSec) implements PaymentResult {}

    static String describe(PaymentResult res) {
        return switch (res) {
            case Approved(var id, var cents) when cents >= 100_00 -> "High-value approved: " + id;
            case Approved(var id, var cents)                      -> "Approved: " + id + " (" + cents + ")";
            case Declined(var reason, var wait)                   -> "Declined (" + reason + "), wait " + wait + "s";
        };
    }
}
\end{lstlisting}
En el sistema de tipos genéricos de Java (implementado con borrado de tipos a \texttt{Object} o a su cota superior para mantener compatibilidad binaria), la varianza en sitios de uso se rige por la regla **PECS** (*Producer Extends, Consumer Super*, formulada por Joshua Bloch): usa \texttt{List<? extends T>} cuando solo leas elementos \texttt{T} de la estructura productora, y \texttt{Comparator<? super T>} o \texttt{Collection<? super T>} cuando insertes o consumas elementos \texttt{T}.
% studium:end-record eyJpZCI6ICJQQVItMDAxNSIsICJzZWN0aW9uIjogImNhcDAzLWphdmEtanZtLWJ5dGVjb2RlLW1lbW9yaWEtZ2VuZXJpY3Mtc3RyZWFtcyJ9

\begin{tcolorbox}[title={Problema resuelto}, studiumworked, breakable]
% studium:record eyJpZCI6ICJQUkItMDAxMyIsICJzZWN0aW9uIjogImNhcDAzLWphdmEtanZtLWJ5dGVjb2RlLW1lbW9yaWEtZ2VuZXJpY3Mtc3RyZWFtcyJ9

\noindent\textbf{Enunciado}
En la JVM HotSpot de 64 bits con punteros comprimidos activos (\texttt{-XX:+UseCompressedOops}), un objeto \texttt{Double} encajonado (*boxed*) ocupa \(24\text{ B}\) (12 B de cabecera + 8 B del \texttt{double} + 4 B de relleno a múltiplo de 8), más \(4\text{ B}\) de la referencia en un \texttt{ArrayList<Double>} (\(28\text{ B}\) por número), mientras que un array primitivo \texttt{double[]} o un \texttt{DoubleStream} ocupa únicamente \(8\text{ B}\) por número. Calcula cuántos megabytes (dividiendo entre \$\(10^6\)\$) se ahorran al almacenar \(N = 2000000\) de valores en \texttt{double[]} frente a \texttt{ArrayList<Double>}.

\noindent\textbf{Datos}
\begin{itemize}
\item $N = 2000000$ — Dos millones de números en coma flotante
\item $S_{boxed} = 28\,\mathrm{B}$ — 24 B del objeto Double + 4 B de referencia comprimida
\item $S_{prim} = 8\,\mathrm{B}$ — 8 B por elemento en double[]
\end{itemize}

\noindent\textbf{Incógnita}
Magnitud objetivo solicitada en el enunciado

\noindent\textbf{Modelo físico}

\noindent\textit{Sobrecoste de memoria por boxing en la JVM frente a tipos primitivos}

\[
\Delta M_{\text{MB}} = \frac{2 \cdot 10^6 \cdot (28 - 8)}{10^6}
\]

\noindent\textbf{Hipótesis}
\begin{itemize}
\item Evaluación determinista bajo el modelo formal del capítulo
\end{itemize}

\noindent\textbf{Desarrollo}
\begin{enumerate}
\item
Multiplicamos los 2 millones de elementos por los 20 B ahorrados y dividimos entre \(10^6\)

\[
\Delta M_{\text{MB}} = \frac{2 \cdot 10^6 \cdot (28 - 8)}{10^6}
\]

\[
\frac{\left(2000000\cdot \left(28-8\right)\right)}{1000000}
\]
\end{enumerate}

\noindent\textbf{Resultado}

\[
Resultado = 40
\]

\noindent\textbf{Interpretación}
Usar arrays o flujos primitivos ahorra exactamente 40 MB de montículo y elimina 2 millones de objetos del recolector de basura.

\noindent\textbf{Comprobación de los pasos}
\begin{itemize}
\item Paso 1: $40$ — reproducido
\end{itemize}

\noindent\textbf{Comprobación de la respuesta} la aritmética del último paso reproduce la respuesta enunciada a la precisión con que está escrita. Es una reproducción aritmética; no valida el modelo físico ni sus hipótesis.
% studium:end-record eyJpZCI6ICJQUkItMDAxMyIsICJzZWN0aW9uIjogImNhcDAzLWphdmEtanZtLWJ5dGVjb2RlLW1lbW9yaWEtZ2VuZXJpY3Mtc3RyZWFtcyJ9
\end{tcolorbox}

\begin{tcolorbox}[title={Ejercicios}, studiumworked, breakable]

\noindent\textbf{1 [PRB-0015] (Básico, Numérico)}
En la JVM HotSpot con Compressed OOPs, la cabecera de un objeto ordinario mide 8 B (\texttt{Mark Word}) más 4 B (\texttt{Klass Pointer}). ¿Cuántos bytes mide la cabecera antes de alinear campos?

\noindent\textbf{2 [PRB-0016] (Básico, Desarrollo por pasos)}
En \texttt{HashMap} de Java 8+, cuando un cubo (*bucket*) alcanza el umbral \texttt{TREEIFY\_THRESHOLD = 8} colisiones (con capacidad de tabla \(\ge 64\)), la lista enlazada se transforma en un árbol rojo-negro. Si hay 256 elementos colisionando en ese mismo cubo, ¿cuál es la altura máxima \(2 \log_2(256)\) a comparar frente a los 256 saltos lineales?

\noindent\textbf{3 [PRB-0017] (Intermedio, Desarrollo por pasos)}
Un \texttt{ConcurrentHashMap} redimensiona su tabla de 1024 a 2048 cubos manteniendo el factor de carga en \(0.75\). ¿Cuántas entradas admite la nueva tabla antes del siguiente redimensionado (\(2048 \times 3 / 4\))?

\noindent\textbf{4 [PRB-0018] (Intermedio, Diseño)}
Si el tamaño de un montículo de la JVM supera los 32 GiB y desactiva Compressed OOPs, cada referencia pasa de 4 B a 8 B. En un grafo en memoria con 50 millones de referencias, ¿cuántos megabytes (\(10^6\text{ B}\)) adicionales consumen las referencias?
\end{tcolorbox}

\chapter{Monografía de Java (II \(-\) Bloques E\(-\)H): Java Memory Model, Virtual Threads (Loom), Foreign Function \& Memory API y Proyectos}
\label{chap:cap04-java-avanzado-experto-virtual-threads-jmm-spring-proyectos}
% studium:record eyJpZCI6ICJQQVItMDAxOSIsICJzZWN0aW9uIjogImNhcDA0LWphdmEtYXZhbnphZG8tZXhwZXJ0by12aXJ0dWFsLXRocmVhZHMtam1tLXNwcmluZy1wcm95ZWN0b3MifQ==

\noindent\textit{Este capítulo completa la monografía de Java (Bloques E, F, G y H). Analizamos con rigor la innovación más importante de la JVM en la última década: los **Hilos Virtuales** (*Virtual Threads*, JEP 444 en Java 21 LTS), que reconcilian el modelo sencillo de un hilo por petición (*thread-per-request*) con la escalabilidad masiva de la E/S no bloqueante.}
% studium:end-record eyJpZCI6ICJQQVItMDAxOSIsICJzZWN0aW9uIjogImNhcDA0LWphdmEtYXZhbnphZG8tZXhwZXJ0by12aXJ0dWFsLXRocmVhZHMtam1tLXNwcmluZy1wcm95ZWN0b3MifQ==

\begin{tcolorbox}[title={Consejo}, studiumtip, breakable]
% studium:record eyJpZCI6ICJQQVItMDAyMyIsICJzZWN0aW9uIjogImNhcDA0LWphdmEtYXZhbnphZG8tZXhwZXJ0by12aXJ0dWFsLXRocmVhZHMtam1tLXNwcmluZy1wcm95ZWN0b3MifQ==

No reutilices en un pool (\texttt{ThreadPool}) los hilos virtuales: los hilos virtuales son baratos y efímeros; crea uno nuevo por cada tarea con \texttt{Executors.newVirtualThreadPerTaskExecutor()} y usa un \texttt{Semaphore} cuando necesites limitar el número de llamadas simultáneas a una base de datos externa.
% studium:end-record eyJpZCI6ICJQQVItMDAyMyIsICJzZWN0aW9uIjogImNhcDA0LWphdmEtYXZhbnphZG8tZXhwZXJ0by12aXJ0dWFsLXRocmVhZHMtam1tLXNwcmluZy1wcm95ZWN0b3MifQ==
\end{tcolorbox}

\begin{tcolorbox}[title={Definición}, studiumdef, breakable]
% studium:record eyJpZCI6ICJQQVItMDAyMiIsICJzZWN0aW9uIjogImNhcDA0LWphdmEtYXZhbnphZG8tZXhwZXJ0by12aXJ0dWFsLXRocmVhZHMtam1tLXNwcmluZy1wcm95ZWN0b3MifQ==

Hilo Virtual en la JVM (Virtual Thread, JEP 444): instancia de \texttt{java.lang.Thread} planificada en espacio de usuario por la máquina virtual sobre un conjunto reducido de hilos portadores del sistema operativo (*carrier threads*), cuya pila de marcos crece y se encoge dinámicamente en el montículo durante bloqueos de E/S.
% studium:end-record eyJpZCI6ICJQQVItMDAyMiIsICJzZWN0aW9uIjogImNhcDA0LWphdmEtYXZhbnphZG8tZXhwZXJ0by12aXJ0dWFsLXRocmVhZHMtam1tLXNwcmluZy1wcm95ZWN0b3MifQ==
\end{tcolorbox}

\begin{tcolorbox}[title={Autoficha}, studiumcheck, breakable]
% studium:record eyJpZCI6ICJQQVItMDAyNCIsICJzZWN0aW9uIjogImNhcDA0LWphdmEtYXZhbnphZG8tZXhwZXJ0by12aXJ0dWFsLXRocmVhZHMtam1tLXNwcmluZy1wcm95ZWN0b3MifQ==

¿Qué es el *anclaje de hilo portador* (*carrier thread pinning*) cuando un hilo virtual se bloquea en E/S dentro de un bloque \texttt{synchronized} tradicional en Java 21 y por qué sustituirlo por \texttt{ReentrantLock} evita el anclaje?
% studium:end-record eyJpZCI6ICJQQVItMDAyNCIsICJzZWN0aW9uIjogImNhcDA0LWphdmEtYXZhbnphZG8tZXhwZXJ0by12aXJ0dWFsLXRocmVhZHMtam1tLXNwcmluZy1wcm95ZWN0b3MifQ==
\end{tcolorbox}

\section{Bloques E y F: Java Memory Model, Virtual Threads (Project Loom), StructuredTaskScope y API Panama (FFM)}
% studium:record eyJpZCI6ICJQQVItMDAyMCIsICJzZWN0aW9uIjogImNhcDA0LWphdmEtYXZhbnphZG8tZXhwZXJ0by12aXJ0dWFsLXRocmVhZHMtam1tLXNwcmluZy1wcm95ZWN0b3MifQ==

El **Java Memory Model (JMM)** define formalmente cuándo las escrituras de un hilo son visibles para otro hilo a través de la relación transitiva *happens-before*: la liberación de un monitor (\texttt{synchronized} o \texttt{Lock.unlock()}) *happens-before* toda adquisición posterior del mismo monitor, y una escritura en un campo \texttt{volatile} *happens-before* toda lectura posterior de ese campo. En Java 21 LTS, **Project Loom** introdujo los **Virtual Threads** (\texttt{Thread.ofVirtual().start(...)} o \texttt{Executors.newVirtualThreadPerTaskExecutor()}): hilos gestionados por la JVM que se montan sobre un pequeño grupo de hilos del sistema operativo (*carrier threads* en un \texttt{ForkJoinPool}). Cuando el código ejecutado dentro de un hilo virtual realiza una operación bloqueante de red o base de datos (como \texttt{Socket.read()} o \texttt{Thread.sleep()}), el runtime de la JVM **desmonta** automáticamente la pila de continuación del hilo virtual hacia el montículo y deja libre al hilo portador del sistema operativo para ejecutar otro hilo virtual:
\begin{lstlisting}[style=studium,language={Java}]
import java.util.concurrent.Executors;
import java.util.stream.IntStream;

public final class VirtualThreadDispatcher {
    public static long runConcurrentTasks(int taskCount) throws Exception {
        try (var executor = Executors.newVirtualThreadPerTaskExecutor()) {
            var futures = IntStream.range(0, taskCount)
                .mapToObj(i -> executor.submit(() -> (long) i * 2L))
                .toList();
            long total = 0L;
            for (var f : futures) { total += f.get(); }
            return total;
        }
    }
}
\end{lstlisting}
Por la Ley de Little (\(L = \lambda \cdot W\)), un servidor ya no está limitado por el número de hilos de plataforma (\(\approx 200\)\(-\)\(1000\)), sino que puede sostener cientos de miles de peticiones bloqueadas en E/S escribiendo código imperativo directo, fácil de depurar y con trazas de pila legibles.
% studium:end-record eyJpZCI6ICJQQVItMDAyMCIsICJzZWN0aW9uIjogImNhcDA0LWphdmEtYXZhbnphZG8tZXhwZXJ0by12aXJ0dWFsLXRocmVhZHMtam1tLXNwcmluZy1wcm95ZWN0b3MifQ==

\section{Bloques G y H: Ecosistema (Maven, Gradle, JUnit 5, Spring Boot, GraalVM) y los 4 proyectos en Java}
% studium:record eyJpZCI6ICJQQVItMDAyMSIsICJzZWN0aW9uIjogImNhcDA0LWphdmEtYXZhbnphZG8tZXhwZXJ0by12aXJ0dWFsLXRocmVhZHMtam1tLXNwcmluZy1wcm95ZWN0b3MifQ==

Asimismo, Java 21/22 reemplaza a JNI y \texttt{sun.misc.Unsafe} con la **Foreign Function \& Memory API (Project Panama)**: \texttt{Arena.ofConfined()} y \texttt{MemorySegment} permiten gestionar memoria nativa fuera del montículo (*off-heap*) con límites espaciales y temporales comprobados determinísticamente. En el Bloque G (Ecosistema), estudiamos la gestión declarativa de proyectos con \texttt{Maven} (\texttt{pom.xml}) y \texttt{Gradle} (\texttt{build.gradle.kts}), pruebas unitarias y de integración con \texttt{JUnit 5} y \texttt{Testcontainers}, microservicios con \texttt{Spring Boot 3} (\texttt{spring.threads.virtual.enabled=true}) y compilación nativa con \texttt{GraalVM Native Image}.
% studium:end-record eyJpZCI6ICJQQVItMDAyMSIsICJzZWN0aW9uIjogImNhcDA0LWphdmEtYXZhbnphZG8tZXhwZXJ0by12aXJ0dWFsLXRocmVhZHMtam1tLXNwcmluZy1wcm95ZWN0b3MifQ==

\begin{tcolorbox}[title={Problema resuelto}, studiumworked, breakable]
% studium:record eyJpZCI6ICJQUkItMDAxOSIsICJzZWN0aW9uIjogImNhcDA0LWphdmEtYXZhbnphZG8tZXhwZXJ0by12aXJ0dWFsLXRocmVhZHMtam1tLXNwcmluZy1wcm95ZWN0b3MifQ==

\noindent\textbf{Enunciado}
Según la Ley de Little (\(\lambda_{\max} = L_{\max} / W\)), un microservicio donde cada petición pasa \(W = 0.25\text{ s}\) esperando respuestas de APIs externas está limitado con hilos tradicionales de plataforma a \(L_{\text{plat}} = 400\) hilos (caudal máximo \(400 / 0.25 = 1600\text{ req/s}\)). Al migrar a **Virtual Threads** en Java 21 soportando \(L_{\text{virt}} = 20000\) hilos virtuales simultáneos con la misma latencia \(W = 0.25\text{ s}\), calcula cuántas peticiones por segundo adicionales puede sostener el servidor (\$20000 / 0.25 - 400 / 0.25\$).

\noindent\textbf{Datos}
\begin{itemize}
\item $L_{plat} = 400\,\mathrm{hilos}$ — Pool clásico de hilos del sistema operativo
\item $L_{virt} = 20000\,\mathrm{hilos}$ — Hilos virtuales simultáneos en Java 21
\item $W = 0.25\,\mathrm{s}$ — Tiempo medio de espera de E/S por petición
\end{itemize}

\noindent\textbf{Incógnita}
Magnitud objetivo solicitada en el enunciado

\noindent\textbf{Modelo físico}

\noindent\textit{Incremento de caudal máximo por Ley de Little con Virtual Threads (JEP 444)}

\[
\Delta \lambda = \frac{20000}{0.25} - \frac{400}{0.25}
\]

\noindent\textbf{Hipótesis}
\begin{itemize}
\item Evaluación determinista bajo el modelo formal del capítulo
\end{itemize}

\noindent\textbf{Desarrollo}
\begin{enumerate}
\item
Restamos el caudal antiguo (1600 req/s) al nuevo caudal (80000 req/s)

\[
\Delta \lambda = \frac{20000}{0.25} - \frac{400}{0.25}
\]

\[
\left(\frac{20000}{0.25}-\frac{400}{0.25}\right)
\]
\end{enumerate}

\noindent\textbf{Resultado}

\[
Resultado = 78400
\]

\noindent\textbf{Interpretación}
El servidor incrementa su capacidad en 78 400 peticiones/s (pasando de 1600 req/s a 80 000 req/s).

\noindent\textbf{Comprobación de los pasos}
\begin{itemize}
\item Paso 1: $78400$ — reproducido
\end{itemize}

\noindent\textbf{Comprobación de la respuesta} la aritmética del último paso reproduce la respuesta enunciada a la precisión con que está escrita. Es una reproducción aritmética; no valida el modelo físico ni sus hipótesis.
% studium:end-record eyJpZCI6ICJQUkItMDAxOSIsICJzZWN0aW9uIjogImNhcDA0LWphdmEtYXZhbnphZG8tZXhwZXJ0by12aXJ0dWFsLXRocmVhZHMtam1tLXNwcmluZy1wcm95ZWN0b3MifQ==
\end{tcolorbox}

\begin{tcolorbox}[title={Ejercicios}, studiumworked, breakable]

\noindent\textbf{1 [PRB-0021] (Básico, Numérico)}
Por la Ley de Little (\(L = \lambda \cdot W\)), si un sistema recibe \(\lambda = 2500\) peticiones/s y cada una tarda \(W = 0.2\) segundos, ¿cuántas peticiones hay en vuelo simultáneamente en promedio?

\noindent\textbf{2 [PRB-0022] (Básico, Numérico)}
En una CPU de 16 núcleos lógicos, el \texttt{ForkJoinPool} portador de hilos virtuales crea por defecto un hilo portador (*carrier thread*) por núcleo. ¿Cuántos hilos del sistema operativo mantienen activos los hilos virtuales?

\noindent\textbf{3 [PRB-0023] (Intermedio, Desarrollo por pasos)}
En el código \texttt{VirtualThreadDispatcher.runConcurrentTasks(100)}, se suma \texttt{i * 2} para \(i = 0, 1, \dots, 99\). ¿Cuál es el valor exacto de la suma (\(99 \times 100\))?

\noindent\textbf{4 [PRB-0024] (Intermedio, Diseño)}
Una \texttt{Arena.ofConfined()} de Project Panama reserva 64 segmentos nativos de 4096 bytes cada uno. ¿Cuántos bytes de memoria *off-heap* se liberan determinísticamente al cerrarse el bloque \texttt{try-with-resources}?
\end{tcolorbox}

\part{Monografías de Sistemas Cloud y Móviles: Go, Kotlin y Swift}

\chapter{Monografía de Go (I \(-\) Bloques A\(-\)D): Filosofía, Compilación Estática, Slices, Maps, Interfaces Implícitas y Manejo de Errores}
\label{chap:cap05-go-fundamentos-slices-interfaces-errores-genericos}
% studium:record eyJpZCI6ICJQQVItMDAyNSIsICJzZWN0aW9uIjogImNhcDA1LWdvLWZ1bmRhbWVudG9zLXNsaWNlcy1pbnRlcmZhY2VzLWVycm9yZXMtZ2VuZXJpY29zIn0=

\noindent\textit{Este capítulo inicia la monografía autónoma de **Go (Golang 1.22+)** (Bloques A, B, C y D). Go es el lenguaje fundacional de la infraestructura nativa en la nube (Docker, Kubernetes, etcd, Prometheus, Terraform, Caddy). Estudiaremos por qué su deliberada economía de conceptos produce sistemas distribuidos mantenibles por grandes equipos.}
% studium:end-record eyJpZCI6ICJQQVItMDAyNSIsICJzZWN0aW9uIjogImNhcDA1LWdvLWZ1bmRhbWVudG9zLXNsaWNlcy1pbnRlcmZhY2VzLWVycm9yZXMtZ2VuZXJpY29zIn0=

\begin{tcolorbox}[title={Consejo}, studiumtip, breakable]
% studium:record eyJpZCI6ICJQQVItMDAyOSIsICJzZWN0aW9uIjogImNhcDA1LWdvLWZ1bmRhbWVudG9zLXNsaWNlcy1pbnRlcmZhY2VzLWVycm9yZXMtZ2VuZXJpY29zIn0=

Sigue el axioma idiomático de Go: *«Acepta interfaces (pequeñas, como \texttt{io.Reader}), devuelve estructuras concretas»*, y evalúa siempre los argumentos de \texttt{defer} teniendo en cuenta que se capturan en el momento en que se ejecuta la sentencia \texttt{defer}, no cuando retorna la función.
% studium:end-record eyJpZCI6ICJQQVItMDAyOSIsICJzZWN0aW9uIjogImNhcDA1LWdvLWZ1bmRhbWVudG9zLXNsaWNlcy1pbnRlcmZhY2VzLWVycm9yZXMtZ2VuZXJpY29zIn0=
\end{tcolorbox}

\begin{tcolorbox}[title={Definición}, studiumdef, breakable]
% studium:record eyJpZCI6ICJQQVItMDAyOCIsICJzZWN0aW9uIjogImNhcDA1LWdvLWZ1bmRhbWVudG9zLXNsaWNlcy1pbnRlcmZhY2VzLWVycm9yZXMtZ2VuZXJpY29zIn0=

Expresión de rebanada completa (\texttt{s[i:j:k]} en Go): operación de sub-slicing que produce una nueva rebanada con longitud \texttt{j - i} y capacidad restringida a \texttt{k - i}, forzando que cualquier \texttt{append} posterior sobre la sub-rebanada aloje un nuevo array en lugar de sobrescribir los elementos a partir del índice \texttt{k}.
% studium:end-record eyJpZCI6ICJQQVItMDAyOCIsICJzZWN0aW9uIjogImNhcDA1LWdvLWZ1bmRhbWVudG9zLXNsaWNlcy1pbnRlcmZhY2VzLWVycm9yZXMtZ2VuZXJpY29zIn0=
\end{tcolorbox}

\begin{tcolorbox}[title={Autoficha}, studiumcheck, breakable]
% studium:record eyJpZCI6ICJQQVItMDAzMCIsICJzZWN0aW9uIjogImNhcDA1LWdvLWZ1bmRhbWVudG9zLXNsaWNlcy1pbnRlcmZhY2VzLWVycm9yZXMtZ2VuZXJpY29zIn0=

¿Por qué una función en Go que declara \texttt{var p *MyCustomError = nil; return p} como valor de retorno de tipo \texttt{error} hace que el llamador observe \texttt{err != nil}?
% studium:end-record eyJpZCI6ICJQQVItMDAzMCIsICJzZWN0aW9uIjogImNhcDA1LWdvLWZ1bmRhbWVudG9zLXNsaWNlcy1pbnRlcmZhY2VzLWVycm9yZXMtZ2VuZXJpY29zIn0=
\end{tcolorbox}

\section{Bloques A\(-\)C: Filosofía de Go, toolchain integrada (`go`), tipos, valores cero, arrays vs slices y maps}
% studium:record eyJpZCI6ICJQQVItMDAyNiIsICJzZWN0aW9uIjogImNhcDA1LWdvLWZ1bmRhbWVudG9zLXNsaWNlcy1pbnRlcmZhY2VzLWVycm9yZXMtZ2VuZXJpY29zIn0=

Diseñado en Google en 2007 por Robert Griesemer, Rob Pike y Ken Thompson, Go nació como respuesta a los tiempos de compilación masivos y la complejidad accidental de C++ y Java en servidores multinúcleo. La herramienta unificada \texttt{go} integra compilación, gestión de módulos (\texttt{go.mod} / \texttt{go.sum} con selección de versión mínima, *Minimal Version Selection*), formateo canónico (\texttt{gofmt}), análisis estático (\texttt{go vet}), pruebas y *fuzzing* nativo (\texttt{go test -fuzz}). En Go, toda variable no inicializada explícitamente recibe su **valor cero** determinista (\texttt{0}, \texttt{false}, \texttt{""}, \texttt{nil}), y los principios de diseño recomiendan que el valor cero de una estructura sea inmediatamente utilizable sin constructor (como \texttt{var mu sync.Mutex} o \texttt{var buf bytes.Buffer}). Un \texttt{slice} \texttt{[]T} es un descriptor de 24 bytes (en 64 bits) formado por un puntero al array subyacente, la longitud \texttt{len} y la capacidad \texttt{cap}. Si dos rebanadas comparten el mismo array subyacente y una ejecuta \texttt{append} sin superar \texttt{cap}, modifica el array compartido; para independizar la capacidad al hacer sub-slicing se utiliza la expresión de tres índices \texttt{s[low:high:max]}.
% studium:end-record eyJpZCI6ICJQQVItMDAyNiIsICJzZWN0aW9uIjogImNhcDA1LWdvLWZ1bmRhbWVudG9zLXNsaWNlcy1pbnRlcmZhY2VzLWVycm9yZXMtZ2VuZXJpY29zIn0=

\section{Bloque D: Métodos, composición de structs, interfaces implícitas, errores como valores y genéricos}
% studium:record eyJpZCI6ICJQQVItMDAyNyIsICJzZWN0aW9uIjogImNhcDA1LWdvLWZ1bmRhbWVudG9zLXNsaWNlcy1pbnRlcmZhY2VzLWVycm9yZXMtZ2VuZXJpY29zIn0=

Go prescinde por completo de la herencia de clases y la sustituye por **composición de estructuras** y **interfaces satisfechas implícitamente**: un tipo satisface una interfaz como \texttt{io.Reader} (\texttt{Read(p []byte) (n int, err error)}) simplemente implementando sus métodos, sin escribir nunca \texttt{implements}. Esto permite definir interfaces pequeñas de uno o dos métodos en el paquete consumidor en lugar de acoplarse al paquete proveedor. Asimismo, Go trata los errores como valores ordinarios devueltos en la última posición de retorno \texttt{(val, error)}, encadenados con \texttt{fmt.Errorf("contexto: \%w", err)} e inspeccionados mediante \texttt{errors.Is} y \texttt{errors.As}:
\begin{lstlisting}[style=studium,language={Go}]
package storage

import (
    "errors"
    "fmt"
    "io"
)

var ErrEmptyPayload = errors.New("empty payload")

func ReadExactHeader(r io.Reader, size int) ([]byte, error) {
    if size <= 0 {
        return nil, ErrEmptyPayload
    }
    buf := make([]byte, size)
    if _, err := io.ReadFull(r, buf); err != nil {
        return nil, fmt.Errorf("read exact header (%d bytes): %w", size, err)
    }
    return buf, nil
}
\end{lstlisting}
Un matiz técnico esencial de las interfaces en Go es su representación interna de dos palabras (\texttt{tab} / tipo dinámico y \texttt{data} / valor dinámico): una interfaz solo es \texttt{== nil} si **ambos** punteros internos son nulos; asignar un puntero tipado \texttt{(*MyError)(nil)} a una variable \texttt{var err error} produce una interfaz **no nula** (\texttt{err != nil}), por lo que las funciones deben devolver siempre el literal \texttt{nil} directamente cuando no hay error.
% studium:end-record eyJpZCI6ICJQQVItMDAyNyIsICJzZWN0aW9uIjogImNhcDA1LWdvLWZ1bmRhbWVudG9zLXNsaWNlcy1pbnRlcmZhY2VzLWVycm9yZXMtZ2VuZXJpY29zIn0=

\begin{tcolorbox}[title={Problema resuelto}, studiumworked, breakable]
% studium:record eyJpZCI6ICJQUkItMDAyNSIsICJzZWN0aW9uIjogImNhcDA1LWdvLWZ1bmRhbWVudG9zLXNsaWNlcy1pbnRlcmZhY2VzLWVycm9yZXMtZ2VuZXJpY29zIn0=

\noindent\textbf{Enunciado}
En un programa en Go, se crea un slice \texttt{base := make([]int64, 100, 200)} (cada \texttt{int64} ocupa \(8\text{ B}\)) y se obtiene la sub-rebanada de tres índices \texttt{sub := base[20:60:80]}. Calcula la suma de la longitud \texttt{len(sub)} (\$60 - 20\$) más la capacidad \texttt{cap(sub)} (\$80 - 20\$) multiplicada por los \(8\text{ B}\) de cada elemento.

\noindent\textbf{Datos}
\begin{itemize}
\item $i = 20$ — Índice inicial
\item $j = 60$ — Índice final de longitud
\item $k = 80$ — Índice límite de capacidad
\end{itemize}

\noindent\textbf{Incógnita}
Magnitud objetivo solicitada en el enunciado

\noindent\textbf{Modelo físico}

\noindent\textit{Cálculo de len y cap en expresiones de rebanada de tres índices s[i:j:k] en Go}

\[
R = ((60 - 20) + (80 - 20)) \cdot 8
\]

\noindent\textbf{Hipótesis}
\begin{itemize}
\item Evaluación determinista bajo el modelo formal del capítulo
\end{itemize}

\noindent\textbf{Desarrollo}
\begin{enumerate}
\item
Sumamos len (40) + cap (60) = 100 y multiplicamos por 8 bytes

\[
R = ((60 - 20) + (80 - 20)) \cdot 8
\]

\[
\left(\left(\left(60-20\right)+\left(80-20\right)\right)\cdot 8\right)
\]
\end{enumerate}

\noindent\textbf{Resultado}

\[
Resultado = 800
\]

\noindent\textbf{Interpretación}
La longitud de sub es 40, su capacidad es 60 y la métrica ponderada por 8 bytes es 800.

\noindent\textbf{Comprobación de los pasos}
\begin{itemize}
\item Paso 1: $800$ — reproducido
\end{itemize}

\noindent\textbf{Comprobación de la respuesta} la aritmética del último paso reproduce la respuesta enunciada a la precisión con que está escrita. Es una reproducción aritmética; no valida el modelo físico ni sus hipótesis.
% studium:end-record eyJpZCI6ICJQUkItMDAyNSIsICJzZWN0aW9uIjogImNhcDA1LWdvLWZ1bmRhbWVudG9zLXNsaWNlcy1pbnRlcmZhY2VzLWVycm9yZXMtZ2VuZXJpY29zIn0=
\end{tcolorbox}

\begin{tcolorbox}[title={Ejercicios}, studiumworked, breakable]

\noindent\textbf{1 [PRB-0027] (Básico, Numérico)}
En Go de 64 bits, el encabezado de un slice (\texttt{SliceHeader}) consta de \texttt{Data} (8 B), \texttt{Len} (8 B) y \texttt{Cap} (8 B). ¿Cuántos bytes ocupa en la pila?

\noindent\textbf{2 [PRB-0028] (Básico, Numérico)}
En Go, cada cubo (*bucket* \texttt{bmap}) de un \texttt{map[K]V} almacena exactamente 8 pares clave-valor. Si un mapa tiene \(2^6 = 64\) cubos, ¿cuántos pares clave-valor caben en esos 64 cubos sin cubos de desbordamiento (\(64 \times 8\))?

\noindent\textbf{3 [PRB-0029] (Intermedio, Desarrollo por pasos)}
Para la rebanada \texttt{s[15:45:50]}, ¿cuántos elementos pueden añadirse con \texttt{append} antes de superar la capacidad y forzar la reserva de un nuevo array (\((50 - 15) - (45 - 15)\))?

\noindent\textbf{4 [PRB-0030] (Intermedio, Diseño)}
Una interfaz \texttt{io.Reader} en Go se representa como un par de palabras \texttt{(itab*, data*)} de 8 bytes cada una. ¿Cuántos bytes ocupan 500 descriptores de interfaz en un slice \texttt{[]io.Reader}?
\end{tcolorbox}

\chapter{Monografía de Go (II \(-\) Bloques E\(-\)H): Planificador G-M-P, Goroutines, Canales, Context, Perfilado pprof y Proyectos}
\label{chap:cap06-go-avanzado-experto-goroutines-gmp-canales-context-proyectos}
% studium:record eyJpZCI6ICJQQVItMDAzMSIsICJzZWN0aW9uIjogImNhcDA2LWdvLWF2YW56YWRvLWV4cGVydG8tZ29yb3V0aW5lcy1nbXAtY2FuYWxlcy1jb250ZXh0LXByb3llY3RvcyJ9

\noindent\textit{Este capítulo completa la monografía de Go (Bloques E, F, G y H). Analizamos el motor de ejecución de Go: su planificador de **goroutines G-M-P**, los patrones de canales y \texttt{context.Context} para evitar fugas de goroutines, el recolector de basura tricolor concurrente de baja latencia y las herramientas integradas de diagnóstico de producción (\texttt{pprof}, \texttt{trace}, \texttt{-race}).}
% studium:end-record eyJpZCI6ICJQQVItMDAzMSIsICJzZWN0aW9uIjogImNhcDA2LWdvLWF2YW56YWRvLWV4cGVydG8tZ29yb3V0aW5lcy1nbXAtY2FuYWxlcy1jb250ZXh0LXByb3llY3RvcyJ9
% studium:record eyJpZCI6ICJGSUctMDAwMSIsICJzZWN0aW9uIjogImNhcDA2LWdvLWF2YW56YWRvLWV4cGVydG8tZ29yb3V0aW5lcy1nbXAtY2FuYWxlcy1jb250ZXh0LXByb3llY3RvcyJ9

\par\vspace{\baselineskip}
\noindent\begin{minipage}{\linewidth}
\centering
\begingroup
\setbox0=\hbox{\begin{tikzpicture}
\tikzset{node distance=1.6cm, >=stealth, thick, font=\small}
\node[draw=blue!70!black, fill=blue!8, rounded corners, minimum width=2.8cm, minimum height=0.9cm, align=center] (m1) {Hilo OS $M_1$};
\node[draw=teal!70!black, fill=teal!10, rounded corners, minimum width=2.8cm, minimum height=0.9cm, align=center, below=0.7cm of m1] (p1) {Procesador $P_1$};
\node[draw=orange!75!black, fill=orange!12, rounded corners, minimum width=3.4cm, minimum height=0.9cm, align=center, below=0.7cm of p1] (q1) {Cola Local: $G_1, G_2, G_3$};
\node[draw=blue!70!black, fill=blue!8, rounded corners, minimum width=2.8cm, minimum height=0.9cm, align=center, right=2.2cm of m1] (m2) {Hilo OS $M_2$};
\node[draw=teal!70!black, fill=teal!10, rounded corners, minimum width=2.8cm, minimum height=0.9cm, align=center, below=0.7cm of m2] (p2) {Procesador $P_2$};
\node[draw=orange!75!black, fill=orange!12, rounded corners, minimum width=3.4cm, minimum height=0.9cm, align=center, below=0.7cm of p2] (q2) {Cola Local Vacía};
\draw[->] (m1) -- (p1);
\draw[->] (p1) -- (q1);
\draw[->] (m2) -- (p2);
\draw[->] (p2) -- (q2);
\draw[->, dashed, red!75!black] (q1.east) -- node[above, font=\footnotesize] {Work-Stealing ($50\%$)} (q2.west);
\end{tikzpicture}}%
\count0=\ht0
\advance\count0 by \dp0
\count2=\wd0
\count4=\dimexpr\textheight/100*38\relax
\count6=\linewidth
\ifnum\count0<1
\count0=1
\fi
\ifnum\count2<1
\count2=1
\fi
\divide\count0 by 4096
\divide\count2 by 4096
\divide\count4 by 4096
\divide\count6 by 4096
\ifnum\count0<1
\count0=1
\fi
\ifnum\count2<1
\count2=1
\fi
\ifnum\count4<1
\count4=1
\fi
\ifnum\count6<1
\count6=1
\fi
\count8=\count2
\ifnum\count0>\count4
\count8=\numexpr\count2*\count4/\count0\relax
\fi
\ifnum\count8>\count6
\resizebox{\linewidth}{!}{\begin{tikzpicture}
\tikzset{node distance=1.6cm, >=stealth, thick, font=\small}
\node[draw=blue!70!black, fill=blue!8, rounded corners, minimum width=2.8cm, minimum height=0.9cm, align=center] (m1) {Hilo OS $M_1$};
\node[draw=teal!70!black, fill=teal!10, rounded corners, minimum width=2.8cm, minimum height=0.9cm, align=center, below=0.7cm of m1] (p1) {Procesador $P_1$};
\node[draw=orange!75!black, fill=orange!12, rounded corners, minimum width=3.4cm, minimum height=0.9cm, align=center, below=0.7cm of p1] (q1) {Cola Local: $G_1, G_2, G_3$};
\node[draw=blue!70!black, fill=blue!8, rounded corners, minimum width=2.8cm, minimum height=0.9cm, align=center, right=2.2cm of m1] (m2) {Hilo OS $M_2$};
\node[draw=teal!70!black, fill=teal!10, rounded corners, minimum width=2.8cm, minimum height=0.9cm, align=center, below=0.7cm of m2] (p2) {Procesador $P_2$};
\node[draw=orange!75!black, fill=orange!12, rounded corners, minimum width=3.4cm, minimum height=0.9cm, align=center, below=0.7cm of p2] (q2) {Cola Local Vacía};
\draw[->] (m1) -- (p1);
\draw[->] (p1) -- (q1);
\draw[->] (m2) -- (p2);
\draw[->] (p2) -- (q2);
\draw[->, dashed, red!75!black] (q1.east) -- node[above, font=\footnotesize] {Work-Stealing ($50\%$)} (q2.west);
\end{tikzpicture}}%
\else
\ifnum\count0>\count4
\resizebox{!}{\dimexpr\textheight/100*38\relax}{\begin{tikzpicture}
\tikzset{node distance=1.6cm, >=stealth, thick, font=\small}
\node[draw=blue!70!black, fill=blue!8, rounded corners, minimum width=2.8cm, minimum height=0.9cm, align=center] (m1) {Hilo OS $M_1$};
\node[draw=teal!70!black, fill=teal!10, rounded corners, minimum width=2.8cm, minimum height=0.9cm, align=center, below=0.7cm of m1] (p1) {Procesador $P_1$};
\node[draw=orange!75!black, fill=orange!12, rounded corners, minimum width=3.4cm, minimum height=0.9cm, align=center, below=0.7cm of p1] (q1) {Cola Local: $G_1, G_2, G_3$};
\node[draw=blue!70!black, fill=blue!8, rounded corners, minimum width=2.8cm, minimum height=0.9cm, align=center, right=2.2cm of m1] (m2) {Hilo OS $M_2$};
\node[draw=teal!70!black, fill=teal!10, rounded corners, minimum width=2.8cm, minimum height=0.9cm, align=center, below=0.7cm of m2] (p2) {Procesador $P_2$};
\node[draw=orange!75!black, fill=orange!12, rounded corners, minimum width=3.4cm, minimum height=0.9cm, align=center, below=0.7cm of p2] (q2) {Cola Local Vacía};
\draw[->] (m1) -- (p1);
\draw[->] (p1) -- (q1);
\draw[->] (m2) -- (p2);
\draw[->] (p2) -- (q2);
\draw[->, dashed, red!75!black] (q1.east) -- node[above, font=\footnotesize] {Work-Stealing ($50\%$)} (q2.west);
\end{tikzpicture}}%
\else
\begin{tikzpicture}
\tikzset{node distance=1.6cm, >=stealth, thick, font=\small}
\node[draw=blue!70!black, fill=blue!8, rounded corners, minimum width=2.8cm, minimum height=0.9cm, align=center] (m1) {Hilo OS $M_1$};
\node[draw=teal!70!black, fill=teal!10, rounded corners, minimum width=2.8cm, minimum height=0.9cm, align=center, below=0.7cm of m1] (p1) {Procesador $P_1$};
\node[draw=orange!75!black, fill=orange!12, rounded corners, minimum width=3.4cm, minimum height=0.9cm, align=center, below=0.7cm of p1] (q1) {Cola Local: $G_1, G_2, G_3$};
\node[draw=blue!70!black, fill=blue!8, rounded corners, minimum width=2.8cm, minimum height=0.9cm, align=center, right=2.2cm of m1] (m2) {Hilo OS $M_2$};
\node[draw=teal!70!black, fill=teal!10, rounded corners, minimum width=2.8cm, minimum height=0.9cm, align=center, below=0.7cm of m2] (p2) {Procesador $P_2$};
\node[draw=orange!75!black, fill=orange!12, rounded corners, minimum width=3.4cm, minimum height=0.9cm, align=center, below=0.7cm of p2] (q2) {Cola Local Vacía};
\draw[->] (m1) -- (p1);
\draw[->] (p1) -- (q1);
\draw[->] (m2) -- (p2);
\draw[->] (p2) -- (q2);
\draw[->, dashed, red!75!black] (q1.east) -- node[above, font=\footnotesize] {Work-Stealing ($50\%$)} (q2.west);
\end{tikzpicture}%
\fi
\fi
\endgroup
\par\nopagebreak
{\small \refstepcounter{studiumfigure}\label{fig:fig-go-gmp-scheduler}\textbf{Figura~\thestudiumfigure.} Arquitectura del planificador M:N de Go: cada contexto lógico P (GOMAXPROCS) posee una cola local sin cerrojos de Goroutines (G) y se vincula a un hilo del sistema operativo (M) para ejecutar código Go o robar la mitad de la cola de otro P.\par}
\end{minipage}
\par\vspace{\baselineskip}
% studium:end-record eyJpZCI6ICJGSUctMDAwMSIsICJzZWN0aW9uIjogImNhcDA2LWdvLWF2YW56YWRvLWV4cGVydG8tZ29yb3V0aW5lcy1nbXAtY2FuYWxlcy1jb250ZXh0LXByb3llY3RvcyJ9

\begin{tcolorbox}[title={Consejo}, studiumtip, breakable]
% studium:record eyJpZCI6ICJQQVItMDAzNSIsICJzZWN0aW9uIjogImNhcDA2LWdvLWF2YW56YWRvLWV4cGVydG8tZ29yb3V0aW5lcy1nbXAtY2FuYWxlcy1jb250ZXh0LXByb3llY3RvcyJ9

Antes de lanzar cualquier \texttt{go func() \{ ... \}()}, responde siempre a dos preguntas: (1) ¿bajo qué condición exacta terminará esta goroutine (por ejemplo, cancelación de \texttt{ctx.Done()} o cierre del canal de entrada)? y (2) ¿quién esperará a que termine (\texttt{sync.WaitGroup} o \texttt{errgroup.Group}) antes de cerrar el proceso?
% studium:end-record eyJpZCI6ICJQQVItMDAzNSIsICJzZWN0aW9uIjogImNhcDA2LWdvLWF2YW56YWRvLWV4cGVydG8tZ29yb3V0aW5lcy1nbXAtY2FuYWxlcy1jb250ZXh0LXByb3llY3RvcyJ9
\end{tcolorbox}

\begin{tcolorbox}[title={Definición}, studiumdef, breakable]
% studium:record eyJpZCI6ICJQQVItMDAzNCIsICJzZWN0aW9uIjogImNhcDA2LWdvLWF2YW56YWRvLWV4cGVydG8tZ29yb3V0aW5lcy1nbXAtY2FuYWxlcy1jb250ZXh0LXByb3llY3RvcyJ9

Fuga de Goroutines (Goroutine Leak): defecto de recursos en Go en el que una goroutine queda suspendida permanentemente en una operación de envío o recepción de canal que ningún otro participante completará jamás, impidiendo que el recolector de basura libere su pila y las variables referenciadas por ella.
% studium:end-record eyJpZCI6ICJQQVItMDAzNCIsICJzZWN0aW9uIjogImNhcDA2LWdvLWF2YW56YWRvLWV4cGVydG8tZ29yb3V0aW5lcy1nbXAtY2FuYWxlcy1jb250ZXh0LXByb3llY3RvcyJ9
\end{tcolorbox}

\begin{tcolorbox}[title={Autoficha}, studiumcheck, breakable]
% studium:record eyJpZCI6ICJQQVItMDAzNiIsICJzZWN0aW9uIjogImNhcDA2LWdvLWF2YW56YWRvLWV4cGVydG8tZ29yb3V0aW5lcy1nbXAtY2FuYWxlcy1jb250ZXh0LXByb3llY3RvcyJ9

¿Qué ocurre exactamente en tiempo de ejecución en Go cuando una goroutine intenta enviar un valor a un canal cerrado (\texttt{close(ch)}) frente a cuando intenta recibir de un canal cerrado?
% studium:end-record eyJpZCI6ICJQQVItMDAzNiIsICJzZWN0aW9uIjogImNhcDA2LWdvLWF2YW56YWRvLWV4cGVydG8tZ29yb3V0aW5lcy1nbXAtY2FuYWxlcy1jb250ZXh0LXByb3llY3RvcyJ9
\end{tcolorbox}

\section{Bloques E y F: Planificador G-M-P, goroutines, canales, select, context.Context, sync/errgroup y análisis de escape}
% studium:record eyJpZCI6ICJQQVItMDAzMiIsICJzZWN0aW9uIjogImNhcDA2LWdvLWF2YW56YWRvLWV4cGVydG8tZ29yb3V0aW5lcy1nbXAtY2FuYWxlcy1jb250ZXh0LXByb3llY3RvcyJ9

Una **goroutine** (\texttt{go fn()}) es un hilo de ejecución ligero gestionado por el runtime de Go cuya pila inicial mide tan solo **\(2\text{ KiB}\)** (\(2048\text{ B}\)) y crece o decrece dinámicamente según la profundidad de llamadas. Como muestra la Figura \textbackslash{}ref\{fig:fig-go-gmp-scheduler\}, el planificador M:N multiplexa \(N\) goroutines (\texttt{G}) sobre \(M\) hilos del sistema operativo (\texttt{M}) utilizando hasta \texttt{GOMAXPROCS} procesadores lógicos (\texttt{P}). Cada \texttt{P} mantiene una cola local de hasta 256 goroutines ejecutables; cuando la cola de un \texttt{P} se vacía, intenta robar la mitad de las goroutines de la cola de otro \texttt{P} elegido al azar (*work-stealing*) o toma trabajo de la cola global. La comunicación entre goroutines sigue el cálculo **CSP** (*Communicating Sequential Processes*): *«No compartas memoria para comunicarte; comunícate para compartir memoria»*. Para evitar que una goroutine quede bloqueada eternamente intentando enviar en un canal cuyo receptor ya ha abandonado la operación (*goroutine leak*), todo ciclo concurrente debe seleccionar también sobre \texttt{<-ctx.Done()} de un \texttt{context.Context}:
\begin{lstlisting}[style=studium,language={Go}]
package worker

import "context"

func FanOutGenerator(ctx context.Context, items []int, out chan<- int) error {
    for _, v := range items {
        select {
        case <-ctx.Done():
            return ctx.Err()
        case out <- v * 2:
        }
    }
    return nil
}
\end{lstlisting}
A partir de Go 1.22, las variables de bucle \texttt{for i, v := range items} tienen ámbito por iteración en lugar de por bucle completo, eliminando el error histórico de capturar la misma variable \texttt{v} en múltiples clausuras concurrentes.
% studium:end-record eyJpZCI6ICJQQVItMDAzMiIsICJzZWN0aW9uIjogImNhcDA2LWdvLWF2YW56YWRvLWV4cGVydG8tZ29yb3V0aW5lcy1nbXAtY2FuYWxlcy1jb250ZXh0LXByb3llY3RvcyJ9

\section{Bloques G y H: Ecosistema (net/http, grpc-go, sqlc, pprof, -race) y los 4 proyectos en Go}
% studium:record eyJpZCI6ICJQQVItMDAzMyIsICJzZWN0aW9uIjogImNhcDA2LWdvLWF2YW56YWRvLWV4cGVydG8tZ29yb3V0aW5lcy1nbXAtY2FuYWxlcy1jb250ZXh0LXByb3llY3RvcyJ9

En el Nivel Experto (Bloque F), el compilador de Go decide automáticamente si una variable local se aloja en la pila de \(2\text{ KiB}\) de la goroutine o en el montículo gestionado mediante **análisis de escape** (*escape analysis*), inspeccionable con \texttt{go build -gcflags="-m"}. En el Bloque G y Bloque H, la biblioteca estándar de Go es tan completa que permite construir servidores HTTP/2 de producción con enrutamiento por patrones (\texttt{net/http} en Go 1.22+), observabilidad (\texttt{log/slog}, \texttt{runtime/pprof}), clientes SQL con pools de conexiones (\texttt{database/sql}) y pruebas de concurrencia instrumentadas con ThreadSanitizer (\texttt{go test -race}). Los cuatro proyectos del bloque incluyen un proxy inverso con balanceo *round-robin* y *health checks*, un grupo de trabajadores (*worker pool*) con apagado ordenado (\texttt{os/signal} + \texttt{context.WithTimeout}), un limitador de tasa distribuido y un servidor gRPC.
% studium:end-record eyJpZCI6ICJQQVItMDAzMyIsICJzZWN0aW9uIjogImNhcDA2LWdvLWF2YW56YWRvLWV4cGVydG8tZ29yb3V0aW5lcy1nbXAtY2FuYWxlcy1jb250ZXh0LXByb3llY3RvcyJ9

\begin{tcolorbox}[title={Problema resuelto}, studiumworked, breakable]
% studium:record eyJpZCI6ICJQUkItMDAzMSIsICJzZWN0aW9uIjogImNhcDA2LWdvLWF2YW56YWRvLWV4cGVydG8tZ29yb3V0aW5lcy1nbXAtY2FuYWxlcy1jb250ZXh0LXByb3llY3RvcyJ9

\noindent\textbf{Enunciado}
Cada goroutine recién creada en Go 1.22 reserva una pila inicial de \(S_G = 2\text{ KiB}\) (\(2048\text{ B}\)), mientras que un hilo nativo de sistema operativo reserva por defecto una pila de \(S_M = 2048\text{ KiB}\) (\(2\text{ MiB}\)). Si un servidor de telemetría en Go mantiene \(N = 100000\) goroutines activas sobre \(M = 16\) hilos del sistema operativo, calcula cuántos mebibytes totales (dividiendo los KiB entre 1024) ocupan las pilas de las 100 000 goroutines más los 16 hilos OS.

\noindent\textbf{Datos}
\begin{itemize}
\item $N = 100000$ — Número de goroutines concurrentes
\item $S_{G} = 2\,\mathrm{KiB}$ — Pila inicial por goroutine
\item $M = 16$ — Hilos del sistema operativo de 2048 KiB
\end{itemize}

\noindent\textbf{Incógnita}
Magnitud objetivo solicitada en el enunciado

\noindent\textbf{Modelo físico}

\noindent\textit{Consumo de memoria de pilas en el modelo M:N de Go}

\[
M_{\text{MiB}} = \frac{100000 \cdot 2 + 16 \cdot 2048}{1024}
\]

\noindent\textbf{Hipótesis}
\begin{itemize}
\item Evaluación determinista bajo el modelo formal del capítulo
\end{itemize}

\noindent\textbf{Desarrollo}
\begin{enumerate}
\item
Sumamos los KiB de las goroutines (200000 KiB) y de los 16 hilos OS (32768 KiB) y dividimos entre 1024

\[
M_{\text{MiB}} = \frac{100000 \cdot 2 + 16 \cdot 2048}{1024}
\]

\[
\frac{\left(\left(100000\cdot 2\right)+\left(16\cdot 2048\right)\right)}{1024}
\]
\end{enumerate}

\noindent\textbf{Resultado}

\[
Resultado = 227.3125
\]

\noindent\textbf{Interpretación}
Las 100 000 goroutines más los 16 hilos del sistema operativo consumen tan solo 227.3125 MiB de memoria de pila.

\noindent\textbf{Comprobación de los pasos}
\begin{itemize}
\item Paso 1: $227.312$ — reproducido
\end{itemize}

\noindent\textbf{Comprobación de la respuesta} la aritmética del último paso reproduce la respuesta enunciada a la precisión con que está escrita. Es una reproducción aritmética; no valida el modelo físico ni sus hipótesis.
% studium:end-record eyJpZCI6ICJQUkItMDAzMSIsICJzZWN0aW9uIjogImNhcDA2LWdvLWF2YW56YWRvLWV4cGVydG8tZ29yb3V0aW5lcy1nbXAtY2FuYWxlcy1jb250ZXh0LXByb3llY3RvcyJ9
\end{tcolorbox}

\begin{tcolorbox}[title={Ejercicios}, studiumworked, breakable]

\noindent\textbf{1 [PRB-0033] (Básico, Numérico)}
¿Cuántos bytes mide la pila inicial de una goroutine en Go (\(2 \times 1024\))?

\noindent\textbf{2 [PRB-0034] (Básico, Numérico)}
En el planificador de Go, un procesador \(P_1\) tiene 180 goroutines en su cola local y un procesador inactivo \(P_2\) le roba la mitad. ¿Cuántas goroutines roba \(P_2\)?

\noindent\textbf{3 [PRB-0035] (Intermedio, Desarrollo por pasos)}
Un canal con búfer \texttt{ch := make(chan uint64, 512)} almacena elementos de 8 bytes. Si hay 300 elementos encolados en el búfer, ¿cuántos elementos adicionales caben antes de bloquear al emisor?

\noindent\textbf{4 [PRB-0036] (Intermedio, Diseño)}
En el problema anterior, ¿cuántos bytes ocupa el búfer circular completo de 512 elementos \texttt{uint64} (\(512 \times 8\))?
\end{tcolorbox}

\chapter{Monografía de Kotlin (Bloques A\(-\)H): Null-Safety, Data/Sealed Classes, Extensiones, Coroutines, Flow y Multiplatform}
\label{chap:cap07-monografia-kotlin-null-safety-coroutines-flow-kmp}
% studium:record eyJpZCI6ICJQQVItMDAzNyIsICJzZWN0aW9uIjogImNhcDA3LW1vbm9ncmFmaWEta290bGluLW51bGwtc2FmZXR5LWNvcm91dGluZXMtZmxvdy1rbXAifQ==

\noindent\textit{Este capítulo ofrece la monografía completa (Bloques \(A-H)\) de **Kotlin (2.0+)**. Desarrollado por JetBrains, Kotlin es el lenguaje oficial prioritario para el desarrollo en Android, una potencia en microservicios de servidor sobre la JVM (Spring Boot, Ktor) y un lenguaje multiplataforma nativo (**Kotlin Multiplatform**, KMP).}
% studium:end-record eyJpZCI6ICJQQVItMDAzNyIsICJzZWN0aW9uIjogImNhcDA3LW1vbm9ncmFmaWEta290bGluLW51bGwtc2FmZXR5LWNvcm91dGluZXMtZmxvdy1rbXAifQ==

\begin{tcolorbox}[title={Consejo}, studiumtip, breakable]
% studium:record eyJpZCI6ICJQQVItMDA0MSIsICJzZWN0aW9uIjogImNhcDA3LW1vbm9ncmFmaWEta290bGluLW51bGwtc2FmZXR5LWNvcm91dGluZXMtZmxvdy1rbXAifQ==

Cuando captures excepciones genéricas (\texttt{catch (e: Throwable)} o \texttt{catch (e: Exception)}) dentro de una corrutina de Kotlin, **relanza siempre \texttt{CancellationException}** (\texttt{if (e is CancellationException) throw e}): tragarse \texttt{CancellationException} rompe el mecanismo de cancelación cooperativa de la concurrencia estructurada.
% studium:end-record eyJpZCI6ICJQQVItMDA0MSIsICJzZWN0aW9uIjogImNhcDA3LW1vbm9ncmFmaWEta290bGluLW51bGwtc2FmZXR5LWNvcm91dGluZXMtZmxvdy1rbXAifQ==
\end{tcolorbox}

\begin{tcolorbox}[title={Definición}, studiumdef, breakable]
% studium:record eyJpZCI6ICJQQVItMDA0MCIsICJzZWN0aW9uIjogImNhcDA3LW1vbm9ncmFmaWEta290bGluLW51bGwtc2FmZXR5LWNvcm91dGluZXMtZmxvdy1rbXAifQ==

Transformación CPS de funciones \texttt{suspend} en Kotlin: reescritura realizada por el compilador en la que toda función \texttt{suspend fun f(a: A): R} se compila en bytecode a \texttt{Object f(A a, Continuation<? super R> completion)}, donde \texttt{completion} almacena el índice \texttt{label} del estado actual y las variables locales vivas entre puntos de suspensión.
% studium:end-record eyJpZCI6ICJQQVItMDA0MCIsICJzZWN0aW9uIjogImNhcDA3LW1vbm9ncmFmaWEta290bGluLW51bGwtc2FmZXR5LWNvcm91dGluZXMtZmxvdy1rbXAifQ==
\end{tcolorbox}

\begin{tcolorbox}[title={Autoficha}, studiumcheck, breakable]
% studium:record eyJpZCI6ICJQQVItMDA0MiIsICJzZWN0aW9uIjogImNhcDA3LW1vbm9ncmFmaWEta290bGluLW51bGwtc2FmZXR5LWNvcm91dGluZXMtZmxvdy1rbXAifQ==

¿Qué diferencia de comportamiento ante el fallo de una corrutina hija existe entre \texttt{coroutineScope \{ ... \}} y \texttt{supervisorScope \{ ... \}} en Kotlin Coroutines?
% studium:end-record eyJpZCI6ICJQQVItMDA0MiIsICJzZWN0aW9uIjogImNhcDA3LW1vbm9ncmFmaWEta290bGluLW51bGwtc2FmZXR5LWNvcm91dGluZXMtZmxvdy1rbXAifQ==
\end{tcolorbox}

\section{Bloques A\(-\)D: Filosofía de Kotlin, interoperabilidad JVM/KMP, null-safety, data/sealed classes y extensiones}
% studium:record eyJpZCI6ICJQQVItMDAzOCIsICJzZWN0aW9uIjogImNhcDA3LW1vbm9ncmFmaWEta290bGluLW51bGwtc2FmZXR5LWNvcm91dGluZXMtZmxvdy1rbXAifQ==

Kotlin resuelve las limitaciones históricas de Java manteniendo una interoperabilidad del 100 \% a nivel de bytecode. Su sistema de tipos incorpora nulabilidad en la raíz de la jerarquía: \texttt{String} no puede almacenar \texttt{null}, mientras que \texttt{String?} exige comprobación explícita mediante el operador de llamada segura \texttt{?.}, el operador Elvis \texttt{?:} o *smart casts* tras verificar \texttt{if (x != null)}. Para modelar dominios sin sobrecoste de ejecución, \texttt{@JvmInline value class UserId(val raw: Long)} envuelve un tipo primitivo en compilación pero se emite como un \texttt{long} desnudo en el bytecode de la JVM. Asimismo, dado que en Kotlin las funciones de orden superior se usan intensivamente, la palabra clave \texttt{inline} inserta el cuerpo de la función y el de las lambdas recibidas directamente en el sitio de llamada, eliminando la creación de objetos clausura y permitiendo preservar el tipo genérico en ejecución con \texttt{reified T}.
% studium:end-record eyJpZCI6ICJQQVItMDAzOCIsICJzZWN0aW9uIjogImNhcDA3LW1vbm9ncmFmaWEta290bGluLW51bGwtc2FmZXR5LWNvcm91dGluZXMtZmxvdy1rbXAifQ==

\section{Bloques E\(-\)H: Funciones inline/reified, DSLs, Coroutines (suspend, CPS, Flow), ecosistema (Ktor, Compose) y proyectos}
% studium:record eyJpZCI6ICJQQVItMDAzOSIsICJzZWN0aW9uIjogImNhcDA3LW1vbm9ncmFmaWEta290bGluLW51bGwtc2FmZXR5LWNvcm91dGluZXMtZmxvdy1rbXAifQ==

La concurrencia en Kotlin se fundamenta en **Coroutines** estructuradas a partir de una única palabra clave en el lenguaje: \texttt{suspend}. El compilador K2 transforma cada \texttt{suspend fun fetch(id: Long): Profile} mediante **Continuation-Passing Style (CPS)** en una función que recibe un parámetro adicional oculto \texttt{Continuation<Profile>} implementado como una máquina de estados sin alojar un hilo por espera:
\begin{lstlisting}[style=studium,escapeinside={(*@}{@*)}]
import kotlinx.coroutines.*

data class NodeHealth(val nodeId: Int, val latencyMs: Long)

suspend fun aggregateClusterLatency(nodeIds: List<Int>): Long = coroutineScope {
    val deferreds: List<Deferred<Long>> = nodeIds.map { id ->
        async(Dispatchers.Default) {
            // Simulaci(*@{\'o}@*)n de consulta paralela por nodo:
            id.toLong() * 5L
        }
    }
    deferreds.awaitAll().sum()
}
\end{lstlisting}
Dentro de \texttt{coroutineScope \{ ... \}}, la concurrencia es estrictamente jerárquica: la función no retorna hasta que todas las corrutinas hijas lanzadas con \texttt{launch} o \texttt{async} han terminado, y si una sola falla con una excepción, todas las demás hermanas reciben cancelación cooperativa inmediata. Para secuencias asíncronas de múltiples valores, Kotlin distingue entre \texttt{Flow<T>} (flujo frío perezoso con contrapresión natural por suspensión) y \texttt{StateFlow<T>} / \texttt{SharedFlow<T>} (flujos calientes observables en arquitecturas reactivas como Jetpack Compose).
% studium:end-record eyJpZCI6ICJQQVItMDAzOSIsICJzZWN0aW9uIjogImNhcDA3LW1vbm9ncmFmaWEta290bGluLW51bGwtc2FmZXR5LWNvcm91dGluZXMtZmxvdy1rbXAifQ==

\begin{tcolorbox}[title={Problema resuelto}, studiumworked, breakable]
% studium:record eyJpZCI6ICJQUkItMDAzNyIsICJzZWN0aW9uIjogImNhcDA3LW1vbm9ncmFmaWEta290bGluLW51bGwtc2FmZXR5LWNvcm91dGluZXMtZmxvdy1rbXAifQ==

\noindent\textbf{Enunciado}
En una aplicación en Kotlin sobre la JVM, una función de orden superior que captura una variable local del entorno se invoca en un bucle caliente \(N = 500000\) veces por segundo. Si la función NO está marcada como \texttt{inline}, cada invocación aloja una instancia de clase anónima de \(S_{\text{lambda}} = 24\text{ B}\). Al marcar la función de orden superior con \texttt{inline}, el alojamiento cae a \(0\text{ B}\). Calcula cuántos megabytes por segundo (dividiendo entre \$\(10^6\)\$) deja de alojar el bucle gracias a \texttt{inline}.

\noindent\textbf{Datos}
\begin{itemize}
\item $N = 500000$ — Invocaciones por segundo
\item $S_{\lambda} = 24\,\mathrm{B}$ — Tamaño del objeto clausura capturado
\end{itemize}

\noindent\textbf{Incógnita}
Magnitud objetivo solicitada en el enunciado

\noindent\textbf{Modelo físico}

\noindent\textit{Ahorro de montículo mediante expansión inline de lambdas con captura en Kotlin}

\[
\Delta M_{\text{MB/s}} = \frac{500000 \cdot 24}{10^6}
\]

\noindent\textbf{Hipótesis}
\begin{itemize}
\item Evaluación determinista bajo el modelo formal del capítulo
\end{itemize}

\noindent\textbf{Desarrollo}
\begin{enumerate}
\item
Multiplicamos las 500000 invocaciones por 24 B y dividimos entre \(10^6\)

\[
\Delta M_{\text{MB/s}} = \frac{500000 \cdot 24}{10^6}
\]

\[
\frac{\left(500000\cdot 24\right)}{1000000}
\]
\end{enumerate}

\noindent\textbf{Resultado}

\[
Resultado = 12
\]

\noindent\textbf{Interpretación}
La palabra clave inline ahorra exactamente 12 MB/s de asignaciones efímeras en el montículo.

\noindent\textbf{Comprobación de los pasos}
\begin{itemize}
\item Paso 1: $12$ — reproducido
\end{itemize}

\noindent\textbf{Comprobación de la respuesta} la aritmética del último paso reproduce la respuesta enunciada a la precisión con que está escrita. Es una reproducción aritmética; no valida el modelo físico ni sus hipótesis.
% studium:end-record eyJpZCI6ICJQUkItMDAzNyIsICJzZWN0aW9uIjogImNhcDA3LW1vbm9ncmFmaWEta290bGluLW51bGwtc2FmZXR5LWNvcm91dGluZXMtZmxvdy1rbXAifQ==
\end{tcolorbox}

\begin{tcolorbox}[title={Ejercicios}, studiumworked, breakable]

\noindent\textbf{1 [PRB-0039] (Básico, Numérico)}
En la función \texttt{aggregateClusterLatency} con \texttt{nodeIds = [1, 2, 3, 4, 5]} donde cada nodo devuelve \texttt{id * 5}, ¿cuál es la suma total devuelta (\((1 + 2 + 3 + 4 + 5) \times 5\))?

\noindent\textbf{2 [PRB-0040] (Básico, Numérico)}
Un \texttt{@JvmInline value class Meter(val v: Double)} se almacena en un array primitivo de 1000 elementos. ¿Cuántos bytes ocupan los 1000 elementos \texttt{double} (\(1000 \times 8\))?

\noindent\textbf{3 [PRB-0041] (Intermedio, Desarrollo por pasos)}
Un \texttt{Flow<Int>} emite 60 elementos a razón de 1 elemento cada 10 ms y el colector aplica \texttt{.buffer(10)} tardando 25 ms por elemento. ¿Cuántos milisegundos tarda el colector en procesar los 60 elementos (\(10 + 60 \times 25\))?

\noindent\textbf{4 [PRB-0042] (Intermedio, Diseño)}
Un proyecto Kotlin Multiplatform comparte 14 000 líneas de lógica de dominio en \texttt{commonMain} entre Android e iOS, requiriendo solo 1500 líneas específicas en \texttt{androidMain} y 1500 en \texttt{iosMain}. ¿Cuántas líneas de código duplicado se ahorran respecto a escribir los 14 000 líneas dos veces?
\end{tcolorbox}

\chapter{Monografía de Swift 6 (Bloques A\(-\)H): Semántica de Valor, Copy-on-Write, Protocolos, ARC y Concurrencia con Actores}
\label{chap:cap08-monografia-swift-valor-cow-protocolos-arc-actores}
% studium:record eyJpZCI6ICJQQVItMDA0MyIsICJzZWN0aW9uIjogImNhcDA4LW1vbm9ncmFmaWEtc3dpZnQtdmFsb3ItY293LXByb3RvY29sb3MtYXJjLWFjdG9yZXMifQ==

\noindent\textit{Este capítulo cierra el Volumen VI con la monografía completa (Bloques \(A-H)\) de **Swift (6.0+)**. Creado por Chris Lattner y el equipo de Apple sobre la infraestructura LLVM, Swift es un lenguaje compilado nativamente sin recolector de basura de trazado que combina semántica de valor profunda, conteo automático de referencias (**ARC**) y verificación estática de aislamiento de actores en tiempo de compilación.}
% studium:end-record eyJpZCI6ICJQQVItMDA0MyIsICJzZWN0aW9uIjogImNhcDA4LW1vbm9ncmFmaWEtc3dpZnQtdmFsb3ItY293LXByb3RvY29sb3MtYXJjLWFjdG9yZXMifQ==

\begin{tcolorbox}[title={Consejo}, studiumtip, breakable]
% studium:record eyJpZCI6ICJQQVItMDA0NyIsICJzZWN0aW9uIjogImNhcDA4LW1vbm9ncmFmaWEtc3dpZnQtdmFsb3ItY293LXByb3RvY29sb3MtYXJjLWFjdG9yZXMifQ==

En firmas de funciones en Swift, prefiere siempre \texttt{func render(item: some Drawable)} sobre \texttt{func render(item: any Drawable)} salvo que necesites almacenar elementos heterogéneos en un mismo array \texttt{[any Drawable]}: \texttt{some} permite al compilador especializar la función y evitar el contenedor existencial.
% studium:end-record eyJpZCI6ICJQQVItMDA0NyIsICJzZWN0aW9uIjogImNhcDA4LW1vbm9ncmFmaWEtc3dpZnQtdmFsb3ItY293LXByb3RvY29sb3MtYXJjLWFjdG9yZXMifQ==
\end{tcolorbox}

\begin{tcolorbox}[title={Definición}, studiumdef, breakable]
% studium:record eyJpZCI6ICJQQVItMDA0NiIsICJzZWN0aW9uIjogImNhcDA4LW1vbm9ncmFmaWEtc3dpZnQtdmFsb3ItY293LXByb3RvY29sb3MtYXJjLWFjdG9yZXMifQ==

Aislamiento de Actor (\texttt{actor} en Swift): tipo por referencia que protege su estado mutable encapsulado mediante un ejecutor serie exclusivo, exigiendo estáticamente en tiempo de compilación que cualquier acceso desde fuera del actor sea asíncrono (\texttt{await}) y que los valores transferidos a través de su frontera cumplan \texttt{Sendable}.
% studium:end-record eyJpZCI6ICJQQVItMDA0NiIsICJzZWN0aW9uIjogImNhcDA4LW1vbm9ncmFmaWEtc3dpZnQtdmFsb3ItY293LXByb3RvY29sb3MtYXJjLWFjdG9yZXMifQ==
\end{tcolorbox}

\begin{tcolorbox}[title={Autoficha}, studiumcheck, breakable]
% studium:record eyJpZCI6ICJQQVItMDA0OCIsICJzZWN0aW9uIjogImNhcDA4LW1vbm9ncmFmaWEtc3dpZnQtdmFsb3ItY293LXByb3RvY29sb3MtYXJjLWFjdG9yZXMifQ==

¿Cómo ocurre la *reentrancia de actores* (*actor reentrancy*) en Swift cuando un método de un \texttt{actor} ejecuta un punto de suspensión \texttt{await} intermedio, y por qué deben revalidarse las invariantes del actor después del \texttt{await}?
% studium:end-record eyJpZCI6ICJQQVItMDA0OCIsICJzZWN0aW9uIjogImNhcDA4LW1vbm9ncmFmaWEtc3dpZnQtdmFsb3ItY293LXByb3RvY29sb3MtYXJjLWFjdG9yZXMifQ==
\end{tcolorbox}

\section{Bloques A\(-\)D: Filosofía de Swift, SIL/LLVM, SwiftPM, semántica de valor (struct/enum), Copy-on-Write y protocolos}
% studium:record eyJpZCI6ICJQQVItMDA0NCIsICJzZWN0aW9uIjogImNhcDA4LW1vbm9ncmFmaWEtc3dpZnQtdmFsb3ItY293LXByb3RvY29sb3MtYXJjLWFjdG9yZXMifQ==

A diferencia de Java o Objective-C, en Swift casi todos los tipos fundamentales de la biblioteca estándar \(-\)incluyendo enteros, flotantes, \texttt{String}, \texttt{Array<T>}, \texttt{Dictionary<K, V>} y \texttt{Set<T>}\(-\) son **tipos por valor** (\texttt{struct}), no clases (\texttt{class}). Para que pasar un \texttt{Array<Double>} de un millón de elementos por valor a una función no implique copiar el búfer inmediatamente, las colecciones de Swift implementan **Copy-on-Write (CoW)**: al asignar \texttt{var b = a}, ambas estructuras comparten el mismo búfer interno en el montículo; solo cuando se ejecuta una mutación como \texttt{b.append(1.0)}, \texttt{Array} comprueba en \(O(1)\) mediante \texttt{isKnownUniquelyReferenced} si el contador de referencias del búfer es mayor que 1 y, únicamente en ese caso, duplica el búfer antes de modificarlo. En el diseño con protocolos (\texttt{protocol}), Swift 5.7+/6.0 distingue explícitamente entre \texttt{some Protocol} (tipo opaco conocido estáticamente en compilación, despachado sin coste) y \texttt{any Protocol} (contenedor existencial dinámico de 5 palabras máquina: 3 palabras de búfer de valor en línea más metadatos de tipo y *witness table*).
% studium:end-record eyJpZCI6ICJQQVItMDA0NCIsICJzZWN0aW9uIjogImNhcDA4LW1vbm9ncmFmaWEtc3dpZnQtdmFsb3ItY293LXByb3RvY29sb3MtYXJjLWFjdG9yZXMifQ==

\section{Bloques E\(-\)H: Gestión de memoria ARC (weak/unowned), concurrencia estricta en Swift 6 (actor, Sendable) y proyectos}
% studium:record eyJpZCI6ICJQQVItMDA0NSIsICJzZWN0aW9uIjogImNhcDA4LW1vbm9ncmFmaWEtc3dpZnQtdmFsb3ItY293LXByb3RvY29sb3MtYXJjLWFjdG9yZXMifQ==

Para las clases (\texttt{class}) y clausuras, Swift gestiona la memoria de forma determinista en tiempo de compilación mediante **ARC** (*Automatic Reference Counting*): el compilador inserta llamadas \texttt{swift\_retain} y \texttt{swift\_release} en la representación intermedia SIL. Como el conteo de referencias no recolecta ciclos aislados, todo ciclo entre instancias de clases o clausuras que capturan \texttt{self} debe romperse declarando una de las aristas como \texttt{weak} (opcional \texttt{T?} que se pone automáticamente a \texttt{nil} al destruirse el objeto) o \texttt{unowned} (no opcional, cuando ambas vidas útiles están ligadas). Finalmente, **Swift 6** garantiza en tiempo de compilación la ausencia de carreras de datos mediante el protocolo marcador \texttt{Sendable} y el tipo de referencia **\texttt{actor}**, que serializa el acceso a su estado mutable interno exigiendo \texttt{await} desde cualquier contexto externo:
\begin{lstlisting}[style=studium]
public actor TelemetryVault {
    private var readings: [String: Double] = [:]
    public init() {}
    public func record(sensor: String, value: Double) {
        readings[sensor] = value
    }
    public func current(sensor: String) -> Double? {
        return readings[sensor]
    }
}
\end{lstlisting}
El capítulo concluye con el ecosistema Swift Package Manager (\texttt{Package.swift}), \texttt{SwiftTesting} (\texttt{@Test}, \texttt{\#expect}), \texttt{SwiftUI} / \texttt{Vapor} en servidor Linux y cuatro proyectos completos en Swift 6.
% studium:end-record eyJpZCI6ICJQQVItMDA0NSIsICJzZWN0aW9uIjogImNhcDA4LW1vbm9ncmFmaWEtc3dpZnQtdmFsb3ItY293LXByb3RvY29sb3MtYXJjLWFjdG9yZXMifQ==

\begin{tcolorbox}[title={Problema resuelto}, studiumworked, breakable]
% studium:record eyJpZCI6ICJQUkItMDA0MyIsICJzZWN0aW9uIjogImNhcDA4LW1vbm9ncmFmaWEtc3dpZnQtdmFsb3ItY293LXByb3RvY29sb3MtYXJjLWFjdG9yZXMifQ==

\noindent\textbf{Enunciado}
En Swift sobre una arquitectura de 64 bits (\(W = 8\text{ B}\)), un contenedor existencial \texttt{any P} para un protocolo ocupa \(5 \times 8 = 40\text{ B}\) (24 B de búfer de valor en línea, 8 B de *Value Witness Table* y 8 B de *Protocol Witness Table*), mientras que un valor concreto \texttt{Int64} pasado mediante un genérico especializado \texttt{some P} ocupa únicamente \(8\text{ B}\). Calcula cuántos bytes ahorra un array de \(N = 10000\) elementos al usar el tipo especializado \texttt{[Int64]} frente a un array existencial \texttt{[any P]}.

\noindent\textbf{Datos}
\begin{itemize}
\item $N = 10000$ — Número de elementos en la colección
\item $S_{any} = 40\,\mathrm{B}$ — Tamaño del contenedor existencial de 5 palabras
\item $S_{some} = 8\,\mathrm{B}$ — Tamaño del tipo concreto especializado
\end{itemize}

\noindent\textbf{Incógnita}
Magnitud objetivo solicitada en el enunciado

\noindent\textbf{Modelo físico}

\noindent\textit{Sobrecoste espacial de contenedores existenciales (any P) frente a tipos concretos (some P) en Swift}

\[
\Delta M = 10000 \cdot (40 - 8)
\]

\noindent\textbf{Hipótesis}
\begin{itemize}
\item Evaluación determinista bajo el modelo formal del capítulo
\end{itemize}

\noindent\textbf{Desarrollo}
\begin{enumerate}
\item
Multiplicamos los 10000 elementos por los 32 B ahorrados por elemento

\[
\Delta M = 10000 \cdot (40 - 8)
\]

\[
\left(10000\cdot \left(40-8\right)\right)
\]
\end{enumerate}

\noindent\textbf{Resultado}

\[
Resultado = 320000
\]

\noindent\textbf{Interpretación}
La especialización estática ahorra 320 000 bytes (un 80 \% de reducción) y elimina el despacho indirecto por tabla de testigos.

\noindent\textbf{Comprobación de los pasos}
\begin{itemize}
\item Paso 1: $320000$ — reproducido
\end{itemize}

\noindent\textbf{Comprobación de la respuesta} la aritmética del último paso reproduce la respuesta enunciada a la precisión con que está escrita. Es una reproducción aritmética; no valida el modelo físico ni sus hipótesis.
% studium:end-record eyJpZCI6ICJQUkItMDA0MyIsICJzZWN0aW9uIjogImNhcDA4LW1vbm9ncmFmaWEtc3dpZnQtdmFsb3ItY293LXByb3RvY29sb3MtYXJjLWFjdG9yZXMifQ==
\end{tcolorbox}

\begin{tcolorbox}[title={Ejercicios}, studiumworked, breakable]

\noindent\textbf{1 [PRB-0045] (Básico, Numérico)}
En Swift de 64 bits (\(W = 8\text{ B}\)), el \texttt{Value Buffer} interno de un contenedor existencial tiene 3 palabras. ¿Hasta cuántos bytes puede medir un \texttt{struct} para caber en línea dentro del contenedor existencial sin alojar en el montículo (\(3 \times 8\))?

\noindent\textbf{2 [PRB-0046] (Básico, Numérico)}
Si se copia un \texttt{Array<Int>} de 50 000 enteros (\texttt{var b = a}) y solo se realizan lecturas sobre \texttt{a} y \texttt{b} sin mutar ninguno de los dos, ¿cuántos elementos se duplican en el montículo gracias a Copy-on-Write?

\noindent\textbf{3 [PRB-0047] (Intermedio, Desarrollo por pasos)}
Un contenedor existencial para una composición de \(K = 3\) protocolos (\texttt{any P1 \& P2 \& P3}) ocupa \((3 + 1 + 3) \times 8\) bytes en 64 bits. ¿Cuántos bytes mide en total?

\noindent\textbf{4 [PRB-0048] (Intermedio, Diseño)}
En la cabecera de un objeto \texttt{class} en Swift de 64 bits, se almacenan el puntero \texttt{isa}/metadatos (8 B) y la palabra de contadores ARC (8 B). ¿Cuántos bytes suman las cabeceras de 2000 instancias?
\end{tcolorbox}
\appendix
\chapter{Notación}

\begingroup
\setlength{\LTpre}{0pt}
\setlength{\LTpost}{0pt}
\begin{longtable}{@{}>{\centering\arraybackslash}p{0.17\linewidth}>{\raggedright\arraybackslash}p{0.55\linewidth}>{\raggedright\arraybackslash}p{0.2\linewidth}@{}}
\textbf{Símbolo} & \textbf{Significado} & \textbf{Unidades} \\
\hline
\endfirsthead
\textbf{Símbolo} & \textbf{Significado} & \textbf{Unidades} \\
\hline
\endhead
$\left( G, \  M, \  P\right)$ & Goroutine, Machine (hilo OS) y Processor (contexto lógico) en el runtime de Go & $\mathrm{entidades}$ \\
$\text{ARC}(o)$ & Contador fuerte/débil/unowned gestionado en compilación en Swift & $\mathrm{referencias}$ \\
$\text{Span}\langle T \rangle$ & Vista contigua de memoria alojada en pila (ref struct) sin asignación en el montículo en .NET & $\mathrm{vista}$ \\
\end{longtable}
\endgroup
\chapter{Hoja de fórmulas}

\noindent Esta parte del borrador aún no tiene material.
\chapter{Soluciones}

\section{Monografía de C\# (I \(-\) Bloques A\(-\)D): Arquitectura .NET/CLR, Tipos Valor vs Referencia, Records, Pattern Matching y LINQ}

\noindent\textbf{Problema resuelto [PRB-0001]}

En una aplicación financiera en C\# sobre .NET de 64 bits, se mantiene un búfer de \(N = 100000\) cotizaciones donde cada cotización tiene \(S_{\text{val}} = 16\text{ B}\) de datos (\texttt{long Timestamp}, \texttt{double Price}). Si se modela como \texttt{readonly record struct}, el array \texttt{Quote[]} ocupa \(100000 \times 16\text{ B}\). Si se modelara como \texttt{class}, cada elemento añadiría \(8\text{ B}\) de puntero en el array más \(16\text{ B}\) de cabecera de objeto en el montículo (\(24\text{ B}\) extra por instancia). Calcula cuántos kilobytes (dividiendo entre 1000) ahorra el diseño con \texttt{struct}.

\noindent\textbf{Respuesta}

\[
Resultado = 2400
\]

\noindent\textbf{Comprobación de la respuesta} la aritmética del último paso reproduce la respuesta enunciada a la precisión con que está escrita. Es una reproducción aritmética; no valida el modelo físico ni sus hipótesis.

\noindent\textbf{Ejercicio 1 [PRB-0003]}

En .NET de 64 bits, la cabecera de todo objeto por referencia (\texttt{class}) consta de \texttt{ObjHeader} (8 B) y \texttt{MethodTable*} (8 B). ¿Cuántos bytes suma esta cabecera?

\noindent\textbf{Respuesta} 16

\noindent\textbf{Ejercicio 2 [PRB-0004]}

Un \texttt{readonly record struct Point3D(double X, double Y, double Z)} contiene tres \texttt{double} de 8 bytes. ¿Cuántos bytes ocupa cada instancia?

\noindent\textbf{Respuesta} 24

\noindent\textbf{Ejercicio 3 [PRB-0005]}

Una consulta LINQ encadena \texttt{.Where(...)} y \texttt{.Take(25)} sobre una secuencia perezosa de 10 000 elementos donde se cumple el predicado en 1 de cada 4 elementos. ¿Cuántos elementos evalúa antes de detenerse (\(25 \times 4\))?

\noindent\textbf{Respuesta} 100

\noindent\textbf{Ejercicio 4 [PRB-0006]}

El recolector generacional del CLR divide el montículo de objetos pequeños en 3 generaciones (\texttt{Gen 0}, \texttt{Gen 1}, \texttt{Gen 2}). Si \texttt{Gen 0} tiene un presupuesto de 4 MiB y \texttt{Gen 1} de 2 MiB, ¿cuántos KiB suman las dos generaciones efímeras (\((4 + 2) \times 1024\))?

\noindent\textbf{Respuesta} 6144

\subsection*{Cálculos reproducidos}

\(\frac{\left(100000\cdot 24\right)}{1000} = 2400\)

\section{Monografía de C\# (II \(-\) Bloques E\(-\)H): Span<T>, Memory<T>, Async/Await, ValueTask, ASP.NET Core y Proyectos}

\noindent\textbf{Problema resuelto [PRB-0007]}

Un servicio de alta frecuencia en ASP.NET Core procesa \(R = 50000\) llamadas por segundo donde el \(90\%\) (\(45000\) llamadas/s) se resuelven síncronamente desde memoria local y solo el \(10\%\) (\(5000\) llamadas/s) requieren E/S asíncrona real. Si cada objeto \texttt{Task<Result>} alojado en el montículo mide \(S_{\text{task}} = 72\text{ B}\), calcula cuántos kilobytes por segundo (dividiendo entre 1000) deja de alojar en \texttt{Gen 0} el servicio al cambiar el tipo de retorno de \texttt{Task<Result>} a \texttt{ValueTask<Result>}.

\noindent\textbf{Respuesta}

\[
Resultado = 3240
\]

\noindent\textbf{Comprobación de la respuesta} la aritmética del último paso reproduce la respuesta enunciada a la precisión con que está escrita. Es una reproducción aritmética; no valida el modelo físico ni sus hipótesis.

\noindent\textbf{Ejercicio 1 [PRB-0009]}

En .NET de 64 bits, un \texttt{Span<byte>} contiene una referencia interior \texttt{ref byte} (8 B) y una longitud \texttt{int} alineada a 8 B. ¿Cuántos bytes ocupa en la pila?

\noindent\textbf{Respuesta} 16

\noindent\textbf{Ejercicio 2 [PRB-0010]}

Si se alquila un búfer con \texttt{ArrayPool<byte>.Shared.Rent(4000)}, el pool devuelve un array cuya longitud es la siguiente potencia de dos (\(4096\)). ¿Cuántos bytes adicionales respecto a los 4000 solicitados tiene el array devuelto?

\noindent\textbf{Respuesta} 96

\noindent\textbf{Ejercicio 3 [PRB-0011]}

Un canal acotado \texttt{Channel.CreateBounded<Job>(new BoundedChannelOptions(128))} contiene actualmente 95 trabajos. ¿Cuántos trabajos más pueden escribirse sin suspender al productor?

\noindent\textbf{Respuesta} 33

\noindent\textbf{Ejercicio 4 [PRB-0012]}

Un contenedor de inyección de dependencias resuelve por cada petición HTTP 1 servicio \texttt{Singleton} (ya creado), 2 servicios \texttt{Scoped} y 4 servicios \texttt{Transient} (invocados 2 veces cada uno en la petición). ¿Cuántas instancias nuevas se crean durante esa petición (\(2 + 4 \times 2\))?

\noindent\textbf{Respuesta} 10

\subsection*{Cálculos reproducidos}

\(\frac{\left(45000\cdot 72\right)}{1000} = 3240\)

\section{Monografía de Java (I \(-\) Bloques A\(-\)D): Arquitectura JVM, Bytecode, Memoria, Records, Sealed Classes, Generics y Streams}

\noindent\textbf{Problema resuelto [PRB-0013]}

En la JVM HotSpot de 64 bits con punteros comprimidos activos (\texttt{-XX:+UseCompressedOops}), un objeto \texttt{Double} encajonado (*boxed*) ocupa \(24\text{ B}\) (12 B de cabecera + 8 B del \texttt{double} + 4 B de relleno a múltiplo de 8), más \(4\text{ B}\) de la referencia en un \texttt{ArrayList<Double>} (\(28\text{ B}\) por número), mientras que un array primitivo \texttt{double[]} o un \texttt{DoubleStream} ocupa únicamente \(8\text{ B}\) por número. Calcula cuántos megabytes (dividiendo entre \$\(10^6\)\$) se ahorran al almacenar \(N = 2000000\) de valores en \texttt{double[]} frente a \texttt{ArrayList<Double>}.

\noindent\textbf{Respuesta}

\[
Resultado = 40
\]

\noindent\textbf{Comprobación de la respuesta} la aritmética del último paso reproduce la respuesta enunciada a la precisión con que está escrita. Es una reproducción aritmética; no valida el modelo físico ni sus hipótesis.

\noindent\textbf{Ejercicio 1 [PRB-0015]}

En la JVM HotSpot con Compressed OOPs, la cabecera de un objeto ordinario mide 8 B (\texttt{Mark Word}) más 4 B (\texttt{Klass Pointer}). ¿Cuántos bytes mide la cabecera antes de alinear campos?

\noindent\textbf{Respuesta} 12

\noindent\textbf{Ejercicio 2 [PRB-0016]}

En \texttt{HashMap} de Java 8+, cuando un cubo (*bucket*) alcanza el umbral \texttt{TREEIFY\_THRESHOLD = 8} colisiones (con capacidad de tabla \(\ge 64\)), la lista enlazada se transforma en un árbol rojo-negro. Si hay 256 elementos colisionando en ese mismo cubo, ¿cuál es la altura máxima \(2 \log_2(256)\) a comparar frente a los 256 saltos lineales?

\noindent\textbf{Respuesta} 16

\noindent\textbf{Ejercicio 3 [PRB-0017]}

Un \texttt{ConcurrentHashMap} redimensiona su tabla de 1024 a 2048 cubos manteniendo el factor de carga en \(0.75\). ¿Cuántas entradas admite la nueva tabla antes del siguiente redimensionado (\(2048 \times 3 / 4\))?

\noindent\textbf{Respuesta} 1536

\noindent\textbf{Ejercicio 4 [PRB-0018]}

Si el tamaño de un montículo de la JVM supera los 32 GiB y desactiva Compressed OOPs, cada referencia pasa de 4 B a 8 B. En un grafo en memoria con 50 millones de referencias, ¿cuántos megabytes (\(10^6\text{ B}\)) adicionales consumen las referencias?

\noindent\textbf{Respuesta} 200

\subsection*{Cálculos reproducidos}

\(\frac{\left(2000000\cdot \left(28-8\right)\right)}{1000000} = 40\)

\section{Monografía de Java (II \(-\) Bloques E\(-\)H): Java Memory Model, Virtual Threads (Loom), Foreign Function \& Memory API y Proyectos}

\noindent\textbf{Problema resuelto [PRB-0019]}

Según la Ley de Little (\(\lambda_{\max} = L_{\max} / W\)), un microservicio donde cada petición pasa \(W = 0.25\text{ s}\) esperando respuestas de APIs externas está limitado con hilos tradicionales de plataforma a \(L_{\text{plat}} = 400\) hilos (caudal máximo \(400 / 0.25 = 1600\text{ req/s}\)). Al migrar a **Virtual Threads** en Java 21 soportando \(L_{\text{virt}} = 20000\) hilos virtuales simultáneos con la misma latencia \(W = 0.25\text{ s}\), calcula cuántas peticiones por segundo adicionales puede sostener el servidor (\$20000 / 0.25 - 400 / 0.25\$).

\noindent\textbf{Respuesta}

\[
Resultado = 78400
\]

\noindent\textbf{Comprobación de la respuesta} la aritmética del último paso reproduce la respuesta enunciada a la precisión con que está escrita. Es una reproducción aritmética; no valida el modelo físico ni sus hipótesis.

\noindent\textbf{Ejercicio 1 [PRB-0021]}

Por la Ley de Little (\(L = \lambda \cdot W\)), si un sistema recibe \(\lambda = 2500\) peticiones/s y cada una tarda \(W = 0.2\) segundos, ¿cuántas peticiones hay en vuelo simultáneamente en promedio?

\noindent\textbf{Respuesta} 500

\noindent\textbf{Ejercicio 2 [PRB-0022]}

En una CPU de 16 núcleos lógicos, el \texttt{ForkJoinPool} portador de hilos virtuales crea por defecto un hilo portador (*carrier thread*) por núcleo. ¿Cuántos hilos del sistema operativo mantienen activos los hilos virtuales?

\noindent\textbf{Respuesta} 16

\noindent\textbf{Ejercicio 3 [PRB-0023]}

En el código \texttt{VirtualThreadDispatcher.runConcurrentTasks(100)}, se suma \texttt{i * 2} para \(i = 0, 1, \dots, 99\). ¿Cuál es el valor exacto de la suma (\(99 \times 100\))?

\noindent\textbf{Respuesta} 9900

\noindent\textbf{Ejercicio 4 [PRB-0024]}

Una \texttt{Arena.ofConfined()} de Project Panama reserva 64 segmentos nativos de 4096 bytes cada uno. ¿Cuántos bytes de memoria *off-heap* se liberan determinísticamente al cerrarse el bloque \texttt{try-with-resources}?

\noindent\textbf{Respuesta} 262144

\subsection*{Cálculos reproducidos}

\(\left(\frac{20000}{0.25}-\frac{400}{0.25}\right) = 78400\)

\section{Monografía de Go (I \(-\) Bloques A\(-\)D): Filosofía, Compilación Estática, Slices, Maps, Interfaces Implícitas y Manejo de Errores}

\noindent\textbf{Problema resuelto [PRB-0025]}

En un programa en Go, se crea un slice \texttt{base := make([]int64, 100, 200)} (cada \texttt{int64} ocupa \(8\text{ B}\)) y se obtiene la sub-rebanada de tres índices \texttt{sub := base[20:60:80]}. Calcula la suma de la longitud \texttt{len(sub)} (\$60 - 20\$) más la capacidad \texttt{cap(sub)} (\$80 - 20\$) multiplicada por los \(8\text{ B}\) de cada elemento.

\noindent\textbf{Respuesta}

\[
Resultado = 800
\]

\noindent\textbf{Comprobación de la respuesta} la aritmética del último paso reproduce la respuesta enunciada a la precisión con que está escrita. Es una reproducción aritmética; no valida el modelo físico ni sus hipótesis.

\noindent\textbf{Ejercicio 1 [PRB-0027]}

En Go de 64 bits, el encabezado de un slice (\texttt{SliceHeader}) consta de \texttt{Data} (8 B), \texttt{Len} (8 B) y \texttt{Cap} (8 B). ¿Cuántos bytes ocupa en la pila?

\noindent\textbf{Respuesta} 24

\noindent\textbf{Ejercicio 2 [PRB-0028]}

En Go, cada cubo (*bucket* \texttt{bmap}) de un \texttt{map[K]V} almacena exactamente 8 pares clave-valor. Si un mapa tiene \(2^6 = 64\) cubos, ¿cuántos pares clave-valor caben en esos 64 cubos sin cubos de desbordamiento (\(64 \times 8\))?

\noindent\textbf{Respuesta} 512

\noindent\textbf{Ejercicio 3 [PRB-0029]}

Para la rebanada \texttt{s[15:45:50]}, ¿cuántos elementos pueden añadirse con \texttt{append} antes de superar la capacidad y forzar la reserva de un nuevo array (\((50 - 15) - (45 - 15)\))?

\noindent\textbf{Respuesta} 5

\noindent\textbf{Ejercicio 4 [PRB-0030]}

Una interfaz \texttt{io.Reader} en Go se representa como un par de palabras \texttt{(itab*, data*)} de 8 bytes cada una. ¿Cuántos bytes ocupan 500 descriptores de interfaz en un slice \texttt{[]io.Reader}?

\noindent\textbf{Respuesta} 8000

\subsection*{Cálculos reproducidos}

\(\left(\left(\left(60-20\right)+\left(80-20\right)\right)\cdot 8\right) = 800\)

\section{Monografía de Go (II \(-\) Bloques E\(-\)H): Planificador G-M-P, Goroutines, Canales, Context, Perfilado pprof y Proyectos}

\noindent\textbf{Problema resuelto [PRB-0031]}

Cada goroutine recién creada en Go 1.22 reserva una pila inicial de \(S_G = 2\text{ KiB}\) (\(2048\text{ B}\)), mientras que un hilo nativo de sistema operativo reserva por defecto una pila de \(S_M = 2048\text{ KiB}\) (\(2\text{ MiB}\)). Si un servidor de telemetría en Go mantiene \(N = 100000\) goroutines activas sobre \(M = 16\) hilos del sistema operativo, calcula cuántos mebibytes totales (dividiendo los KiB entre 1024) ocupan las pilas de las 100 000 goroutines más los 16 hilos OS.

\noindent\textbf{Respuesta}

\[
Resultado = 227.3125
\]

\noindent\textbf{Comprobación de la respuesta} la aritmética del último paso reproduce la respuesta enunciada a la precisión con que está escrita. Es una reproducción aritmética; no valida el modelo físico ni sus hipótesis.

\noindent\textbf{Ejercicio 1 [PRB-0033]}

¿Cuántos bytes mide la pila inicial de una goroutine en Go (\(2 \times 1024\))?

\noindent\textbf{Respuesta} 2048

\noindent\textbf{Ejercicio 2 [PRB-0034]}

En el planificador de Go, un procesador \(P_1\) tiene 180 goroutines en su cola local y un procesador inactivo \(P_2\) le roba la mitad. ¿Cuántas goroutines roba \(P_2\)?

\noindent\textbf{Respuesta} 90

\noindent\textbf{Ejercicio 3 [PRB-0035]}

Un canal con búfer \texttt{ch := make(chan uint64, 512)} almacena elementos de 8 bytes. Si hay 300 elementos encolados en el búfer, ¿cuántos elementos adicionales caben antes de bloquear al emisor?

\noindent\textbf{Respuesta} 212

\noindent\textbf{Ejercicio 4 [PRB-0036]}

En el problema anterior, ¿cuántos bytes ocupa el búfer circular completo de 512 elementos \texttt{uint64} (\(512 \times 8\))?

\noindent\textbf{Respuesta} 4096

\subsection*{Cálculos reproducidos}

\(\frac{\left(\left(100000\cdot 2\right)+\left(16\cdot 2048\right)\right)}{1024} = 227.3125\)

\section{Monografía de Kotlin (Bloques A\(-\)H): Null-Safety, Data/Sealed Classes, Extensiones, Coroutines, Flow y Multiplatform}

\noindent\textbf{Problema resuelto [PRB-0037]}

En una aplicación en Kotlin sobre la JVM, una función de orden superior que captura una variable local del entorno se invoca en un bucle caliente \(N = 500000\) veces por segundo. Si la función NO está marcada como \texttt{inline}, cada invocación aloja una instancia de clase anónima de \(S_{\text{lambda}} = 24\text{ B}\). Al marcar la función de orden superior con \texttt{inline}, el alojamiento cae a \(0\text{ B}\). Calcula cuántos megabytes por segundo (dividiendo entre \$\(10^6\)\$) deja de alojar el bucle gracias a \texttt{inline}.

\noindent\textbf{Respuesta}

\[
Resultado = 12
\]

\noindent\textbf{Comprobación de la respuesta} la aritmética del último paso reproduce la respuesta enunciada a la precisión con que está escrita. Es una reproducción aritmética; no valida el modelo físico ni sus hipótesis.

\noindent\textbf{Ejercicio 1 [PRB-0039]}

En la función \texttt{aggregateClusterLatency} con \texttt{nodeIds = [1, 2, 3, 4, 5]} donde cada nodo devuelve \texttt{id * 5}, ¿cuál es la suma total devuelta (\((1 + 2 + 3 + 4 + 5) \times 5\))?

\noindent\textbf{Respuesta} 75

\noindent\textbf{Ejercicio 2 [PRB-0040]}

Un \texttt{@JvmInline value class Meter(val v: Double)} se almacena en un array primitivo de 1000 elementos. ¿Cuántos bytes ocupan los 1000 elementos \texttt{double} (\(1000 \times 8\))?

\noindent\textbf{Respuesta} 8000

\noindent\textbf{Ejercicio 3 [PRB-0041]}

Un \texttt{Flow<Int>} emite 60 elementos a razón de 1 elemento cada 10 ms y el colector aplica \texttt{.buffer(10)} tardando 25 ms por elemento. ¿Cuántos milisegundos tarda el colector en procesar los 60 elementos (\(10 + 60 \times 25\))?

\noindent\textbf{Respuesta} 1510

\noindent\textbf{Ejercicio 4 [PRB-0042]}

Un proyecto Kotlin Multiplatform comparte 14 000 líneas de lógica de dominio en \texttt{commonMain} entre Android e iOS, requiriendo solo 1500 líneas específicas en \texttt{androidMain} y 1500 en \texttt{iosMain}. ¿Cuántas líneas de código duplicado se ahorran respecto a escribir los 14 000 líneas dos veces?

\noindent\textbf{Respuesta} 14000

\subsection*{Cálculos reproducidos}

\(\frac{\left(500000\cdot 24\right)}{1000000} = 12\)

\section{Monografía de Swift 6 (Bloques A\(-\)H): Semántica de Valor, Copy-on-Write, Protocolos, ARC y Concurrencia con Actores}

\noindent\textbf{Problema resuelto [PRB-0043]}

En Swift sobre una arquitectura de 64 bits (\(W = 8\text{ B}\)), un contenedor existencial \texttt{any P} para un protocolo ocupa \(5 \times 8 = 40\text{ B}\) (24 B de búfer de valor en línea, 8 B de *Value Witness Table* y 8 B de *Protocol Witness Table*), mientras que un valor concreto \texttt{Int64} pasado mediante un genérico especializado \texttt{some P} ocupa únicamente \(8\text{ B}\). Calcula cuántos bytes ahorra un array de \(N = 10000\) elementos al usar el tipo especializado \texttt{[Int64]} frente a un array existencial \texttt{[any P]}.

\noindent\textbf{Respuesta}

\[
Resultado = 320000
\]

\noindent\textbf{Comprobación de la respuesta} la aritmética del último paso reproduce la respuesta enunciada a la precisión con que está escrita. Es una reproducción aritmética; no valida el modelo físico ni sus hipótesis.

\noindent\textbf{Ejercicio 1 [PRB-0045]}

En Swift de 64 bits (\(W = 8\text{ B}\)), el \texttt{Value Buffer} interno de un contenedor existencial tiene 3 palabras. ¿Hasta cuántos bytes puede medir un \texttt{struct} para caber en línea dentro del contenedor existencial sin alojar en el montículo (\(3 \times 8\))?

\noindent\textbf{Respuesta} 24

\noindent\textbf{Ejercicio 2 [PRB-0046]}

Si se copia un \texttt{Array<Int>} de 50 000 enteros (\texttt{var b = a}) y solo se realizan lecturas sobre \texttt{a} y \texttt{b} sin mutar ninguno de los dos, ¿cuántos elementos se duplican en el montículo gracias a Copy-on-Write?

\noindent\textbf{Respuesta} 0

\noindent\textbf{Ejercicio 3 [PRB-0047]}

Un contenedor existencial para una composición de \(K = 3\) protocolos (\texttt{any P1 \& P2 \& P3}) ocupa \((3 + 1 + 3) \times 8\) bytes en 64 bits. ¿Cuántos bytes mide en total?

\noindent\textbf{Respuesta} 56

\noindent\textbf{Ejercicio 4 [PRB-0048]}

En la cabecera de un objeto \texttt{class} en Swift de 64 bits, se almacenan el puntero \texttt{isa}/metadatos (8 B) y la palabra de contadores ARC (8 B). ¿Cuántos bytes suman las cabeceras de 2000 instancias?

\noindent\textbf{Respuesta} 32000

\subsection*{Cálculos reproducidos}

\(\left(10000\cdot \left(40-8\right)\right) = 320000\)
\chapter{Auditoría de fuentes}

\noindent Estado de las fuentes: PENDING. 31 pendientes. Los conflictos almacenados quedan fuera (0). Las fuentes que la guía almacenada no cita quedan fuera. Recuento de no citadas: 0.

\noindent ¶1: dos testimonios.

\noindent ¶2: dos testimonios.

\noindent ¶3: dos testimonios.

\noindent ¶4: dos testimonios.

\noindent ¶5: dos testimonios.

\noindent ¶6: dos testimonios.

\noindent ¶7: dos testimonios.

\noindent ¶8: dos testimonios.

\noindent ¶9: dos testimonios.

\noindent ¶10: dos testimonios.

\noindent ¶11: dos testimonios.

\noindent ¶12: dos testimonios.

\noindent ¶13: dos testimonios.

\noindent ¶14: dos testimonios.

\noindent ¶15: dos testimonios.

\noindent ¶16: dos testimonios.

\noindent ¶17: dos testimonios.

\noindent ¶18: dos testimonios.

\noindent ¶19: dos testimonios.

\noindent ¶20: dos testimonios. Ecuaciones: 1 respaldada por una fuente.

\noindent ¶21: dos testimonios.

\noindent ¶22: dos testimonios.

\noindent ¶23: dos testimonios.

\noindent ¶24: dos testimonios.

\noindent ¶25: dos testimonios.

\noindent ¶26: dos testimonios.

\noindent ¶27: dos testimonios.

\noindent ¶28: dos testimonios.

\noindent ¶29: dos testimonios.

\noindent ¶30: dos testimonios.

\noindent ¶31: dos testimonios.

\noindent ¶32: dos testimonios.

\noindent ¶33: dos testimonios.

\noindent ¶34: dos testimonios.

\noindent ¶35: dos testimonios.

\noindent ¶36: dos testimonios.

\noindent ¶37: dos testimonios.

\noindent ¶38: dos testimonios.

\noindent ¶39: dos testimonios.

\noindent ¶40: dos testimonios.

\noindent ¶41: dos testimonios.

\noindent ¶42: dos testimonios.

\noindent ¶43: dos testimonios.

\noindent ¶44: dos testimonios.

\noindent ¶45: dos testimonios.

\noindent ¶46: dos testimonios.

\noindent ¶47: dos testimonios.

\noindent ¶48: dos testimonios.

\medskip

\noindent Referencia Normativa src-csharp-spec (PENDING) --- Suplemento abierto: --- No es la bibliografía de la guía. --- \P1, \P2, \P3, \P4, \P5, \P6, \P7, \P8, \P9, \P10, \P11, \P12.

\noindent Referencia Normativa src-gamma-gof (PENDING) --- Suplemento abierto: --- No es la bibliografía de la guía. --- \P1, \P2, \P3, \P4, \P5, \P6, \P13, \P14, \P15, \P16, \P17, \P18.

\noindent Referencia Normativa src-kleppmann-ddia (PENDING) --- Suplemento abierto: --- No es la bibliografía de la guía. --- \P7, \P8, \P9, \P10, \P11, \P12, \P19, \P20, \P21, \P22, \P23, \P24, \P25, \P26, \P27, \P28, \P29, \P30, \P31, \P32, \P33, \P34, \P35, \P36.

\noindent Referencia Normativa src-java-jls21 (PENDING) --- Suplemento abierto: --- No es la bibliografía de la guía. --- \P13, \P14, \P15, \P16, \P17, \P18, \P19, \P20, \P21, \P22, \P23, \P24, \P37, \P38, \P39, \P40, \P41, \P42.

\noindent The Go Programming Language Specification \& Memory Model (PENDING) --- Suplemento abierto: --- No es la bibliografía de la guía. --- \P25, \P26, \P27, \P28, \P29, \P30, \P31, \P32, \P33, \P34, \P35, \P36.

\noindent Referencia Normativa src-kotlin-spec (PENDING) --- Suplemento abierto: --- No es la bibliografía de la guía. --- \P37, \P38, \P39, \P40, \P41, \P42.

\noindent The Swift Programming Language (Swift 6.0 Official Book) (PENDING) --- Suplemento abierto: --- No es la bibliografía de la guía. --- \P43, \P44, \P45, \P46, \P47, \P48.

\noindent Referencia Normativa src-pierce-tapl (PENDING) --- Suplemento abierto: --- No es la bibliografía de la guía. --- \P43, \P44, \P45, \P46, \P47, \P48.

\noindent REV-0001: cap01-csharp-dotnet-clr-tipos-valor-referencia-linq scientific

\noindent REV-0002: cap02-csharp-avanzado-experto-span-async-valuetask-aspnet-proyectos scientific

\noindent REV-0003: cap03-java-jvm-bytecode-memoria-generics-streams scientific

\noindent REV-0004: cap04-java-avanzado-experto-virtual-threads-jmm-spring-proyectos scientific

\noindent REV-0005: cap05-go-fundamentos-slices-interfaces-errores-genericos scientific

\noindent REV-0006: cap06-go-avanzado-experto-goroutines-gmp-canales-context-proyectos scientific

\noindent REV-0007: cap07-monografia-kotlin-null-safety-coroutines-flow-kmp scientific

\noindent REV-0008: cap08-monografia-swift-valor-cow-protocolos-arc-actores scientific
\chapter{Plan de estudio}

\section*{Monografías de Plataformas Gestionadas: C\# (.NET 8/9) y Java (JVM 21 LTS)}

\subsection*{\checkmark Monografía de C\# (I \(-\) Bloques A\(-\)D): Arquitectura .NET/CLR, Tipos Valor vs Referencia, Records, Pattern Matching y LINQ}
\begin{itemize}
\item Comprender la arquitectura del Common Language Runtime (CLR), el lenguaje intermedio CIL, la compilación JIT (RyuJIT) y Native AOT en .NET 8/9.
\item Distinguir con precisión semántica y espacial entre tipos por valor (`struct`, `readonly struct`, `ref struct`) y tipos por referencia (`class`, `record`), evitando el boxing involuntario.
\item Dominar la nulabilidad verificada (`string?`), coincidencia de patrones avanzada y consultas declarativas con LINQ (`IEnumerable<T>` vs `IQueryable<T>`).
\end{itemize}
\noindent 2 conceptos planificados.

\subsection*{\checkmark Monografía de C\# (II \(-\) Bloques E\(-\)H): Span<T>, Memory<T>, Async/Await, ValueTask, ASP.NET Core y Proyectos}
\begin{itemize}
\item Programar procesamiento de alto rendimiento sin asignaciones en el montículo usando `Span<T>`, `ReadOnlySpan<T>`, `Memory<T>`, `ArrayPool<T>` y `System.IO.Pipelines`.
\item Comprender la máquina de estados generada por el compilador para `async`/`await`, cuándo utilizar `ValueTask<T>` frente a `Task<T>` y la propagación de `CancellationToken`.
\item Dominar el ecosistema empresarial de .NET (ASP.NET Core, inyección de dependencias, EF Core, Dapper, xUnit, BenchmarkDotNet) y los 4 proyectos integradores.
\end{itemize}
\noindent 2 conceptos planificados.

\subsection*{\checkmark Monografía de Java (I \(-\) Bloques A\(-\)D): Arquitectura JVM, Bytecode, Memoria, Records, Sealed Classes, Generics y Streams}
\begin{itemize}
\item Explicar la arquitectura de la Java Virtual Machine (JVM): subsistema ClassLoader, áreas de memoria (Heap, Metaspace, pilas de hilo), intérprete y compiladores JIT C1/C2.
\item Dominar las innovaciones de Java 17/21 LTS: `record`, `sealed interface`/`permits`, coincidencia de patrones en `switch` y desestructuración de records.
\item Aplicar genéricos con comodines acotados (principio PECS: Producer Extends, Consumer Super), colecciones (`ArrayList`, `HashMap`, `ConcurrentHashMap`) y la API `Stream<T>`.
\end{itemize}
\noindent 2 conceptos planificados.

\subsection*{\checkmark Monografía de Java (II \(-\) Bloques E\(-\)H): Java Memory Model, Virtual Threads (Loom), Foreign Function \& Memory API y Proyectos}
\begin{itemize}
\item Dominar el Java Memory Model (JMM, JLS §17.4): la relación de orden parcial *happens-before*, `volatile`, `VarHandle`, `CompletableFuture` y `ReentrantLock`.
\item Escalar aplicaciones síncronas legibles a millones de tareas concurrentes mediante **Virtual Threads** (JEP 444 en Java 21 LTS) y Concurrencia Estructurada.
\item Integrar el ecosistema empresarial de Java (Maven, Gradle, JUnit 5, Testcontainers, Spring Boot 3, Micrometer/OpenTelemetry) y construir los 4 proyectos del bloque.
\end{itemize}
\noindent 2 conceptos planificados.

\section*{Monografías de Sistemas Cloud y Móviles: Go, Kotlin y Swift}

\subsection*{\checkmark Monografía de Go (I \(-\) Bloques A\(-\)D): Filosofía, Compilación Estática, Slices, Maps, Interfaces Implícitas y Manejo de Errores}
\begin{itemize}
\item Comprender la filosofía de simplicidad ortogonal de Go (Rob Pike, Ken Thompson, Robert Griesemer), su compilador ultrarrápido y el enlazado de binarios estáticos.
\item Dominar la anatomía en memoria de `slice` (puntero, longitud, capacidad) y `map`, el valor cero útil (*zero value*), composición de `struct` por incrustación (*embedding*) y `defer`.
\item Diseñar abstracciones desacopladas mediante interfaces satisfechas implícitamente, manejo explícito de errores como valores (`errors.Is`, `errors.As`, `\%w`) y genéricos de Go 1.18+.
\end{itemize}
\noindent 2 conceptos planificados.

\subsection*{\checkmark Monografía de Go (II \(-\) Bloques E\(-\)H): Planificador G-M-P, Goroutines, Canales, Context, Perfilado pprof y Proyectos}
\begin{itemize}
\item Explicar el funcionamiento interno del planificador M:N con robo de trabajo (*work-stealing*) de Go basado en las entidades G (Goroutine), M (Machine) y P (Processor).
\item Diseñar patrones concurrentes basados en CSP (Tony Hoare) con canales con y sin búfer, `select`, cancelación jerárquica con `context.Context` y `errgroup.Group`.
\item Auditar el escape al montículo (`go build -gcflags='-m'`), detectar carreras de datos (`go test -race`), perfilar CPU/memoria con `pprof` y construir los 4 proyectos en Go.
\end{itemize}
\noindent 2 conceptos planificados.

\subsection*{\checkmark Monografía de Kotlin (Bloques A\(-\)H): Null-Safety, Data/Sealed Classes, Extensiones, Coroutines, Flow y Multiplatform}
\begin{itemize}
\item Dominar el sistema de tipos de Kotlin 2.0 (compilador K2): distinción estricta `T` vs `T?`, smart casts, `data class`, `sealed interface`, `value class` y delegación (`by`).
\item Escribir código expresivo con funciones de orden superior `inline`, tipos genéricos reificados (`reified T`), funciones de extensión y DSLs con receptores (`T.() -> Unit`).
\item Construir sistemas concurrentes y reactivos con **Kotlin Coroutines**: funciones `suspend`, transformaciones CPS, `CoroutineScope`, `Dispatchers`, `Structured Concurrency` y `Flow` / `StateFlow`.
\end{itemize}
\noindent 2 conceptos planificados.

\subsection*{\checkmark Monografía de Swift 6 (Bloques A\(-\)H): Semántica de Valor, Copy-on-Write, Protocolos, ARC y Concurrencia con Actores}
\begin{itemize}
\item Comprender la arquitectura del compilador `swiftc` (SIL: Swift Intermediate Language y LLVM), la semántica de valor (`struct`, `enum`) y la optimización Copy-on-Write (CoW).
\item Diseñar abstracciones con Programación Orientada a Protocolos (`protocol`, tipos asociados, tipos opacos `some P` vs existenciales `any P`).
\item Dominar la gestión determinista de memoria mediante **ARC** (`strong`, `weak`, `unowned`) y el aislamiento estricto de carreras de datos de **Swift 6** (`async`/`await`, `actor`, `@MainActor`, `Sendable`).
\end{itemize}
\noindent 2 conceptos planificados.
\backmatter
\begin{thebibliography}{99}
\bibitem{src1} Alan M. Turing \(-\) On Computable Numbers, with an Application to the Entscheidungsproblem (1936). \url{https://www.cs.virginia.edu/~robins/Turing_Paper_1936.pdf}
\bibitem{src2} John von Neumann \(-\) First Draft of a Report on the EDVAC (1945). \url{https://web.mit.edu/STS.035/www/PDFs/edvac.pdf}
\bibitem{src3} IEEE Standard for Floating-Point Arithmetic (IEEE 754-2019) / Wikipedia Summary. \url{https://en.wikipedia.org/wiki/IEEE_754}
\bibitem{src4} RFC 3629 \(-\) UTF-8, a transformation format of ISO 10646 (IETF Standard). \url{https://datatracker.ietf.org/doc/html/rfc3629}
\bibitem{src5} ACM / IEEE-CS / AAAI Computer Science Curricula 2023 (CS2023). \url{https://csed.acm.org/}
\bibitem{src6} MIT OpenCourseWare 6.006 \(-\) Introduction to Algorithms. \url{https://ocw.mit.edu/courses/6-006-introduction-to-algorithms-spring-2020/}
\bibitem{src7} The Rust Programming Language (Klabnik \& Nichols \(-\) Official Rust Documentation). \url{https://doc.rust-lang.org/book/}
\bibitem{src8} The Rust Reference \(-\) Memory Model, Type System, Lifetimes, and Unsafe Rust. \url{https://doc.rust-lang.org/reference/}
\bibitem{src9} The Rustonomicon \(-\) The Dark Arts of Unsafe Rust. \url{https://doc.rust-lang.org/nomicon/}
\bibitem{src10} C Programming Language Reference (ISO/IEC 9899 C17/C23 \(-\) cppreference). \url{https://en.cppreference.com/w/c}
\bibitem{src11} C++ Reference (ISO/IEC 14882 C++20/C++23 \(-\) cppreference). \url{https://en.cppreference.com/w/cpp}
\bibitem{src12} The Python Language Reference \& Standard Library Documentation (Python 3.12\(-\)3.14). \url{https://docs.python.org/3/reference/index.html}
\bibitem{src13} ECMA-262 \(-\) ECMAScript 2024 Language Specification. \url{https://tc39.es/ecma262/}
\bibitem{src14} The TypeScript Handbook \(-\) Structural Typing, Narrowing, and Type Manipulation. \url{https://www.typescriptlang.org/docs/handbook/intro.html}
\bibitem{src15} Microsoft .NET \& C\# Language Specification (ECMA-334 / .NET 8 \& 9). \url{https://learn.microsoft.com/en-us/dotnet/csharp/}
\bibitem{src16} The Java Language Specification \& JVM Specification (Java SE 21 LTS). \url{https://docs.oracle.com/javase/specs/jls/se21/html/index.html}
\bibitem{src17} The Go Programming Language Specification \& Memory Model. \url{https://go.dev/ref/spec}
\bibitem{src18} Kotlin Language Documentation \& Coroutines Guide (Kotlin 2.0). \url{https://kotlinlang.org/docs/home.html}
\bibitem{src19} The Swift Programming Language (Swift 6.0 Official Book). \url{https://docs.swift.org/swift-book/documentation/the-swift-programming-language/}
\bibitem{src20} E. F. Codd \(-\) A Relational Model of Data for Large Shared Data Banks (1970) \& PostgreSQL Docs. \url{https://www.postgresql.org/docs/current/}
\bibitem{src21} SQLite 3.54 Transactional SQL Database Engine Documentation. \url{https://www.sqlite.org/docs.html}
\bibitem{src22} GNU Bash Reference Manual \& IEEE Std 1003.1 POSIX Shell Specification. \url{https://www.gnu.org/software/bash/manual/bash.html}
\bibitem{src23} WHATWG HTML Living Standard \& W3C CSS Specifications. \url{https://html.spec.whatwg.org/multipage/}
\bibitem{src24} RFC 8259 \(-\) The JavaScript Object Notation (JSON) Data Interchange Format \& YAML 1.2.2. \url{https://datatracker.ietf.org/doc/html/rfc8259}
\bibitem{src25} IEEE Computer Society \(-\) Guide to the Software Engineering Body of Knowledge (SWEBOK v4.0). \url{https://www.computer.org/education/bodies-of-knowledge/software-engineering}
\bibitem{src26} Referencia Normativa src-csharp-spec. \url{https://doc.rust-lang.org/reference/src-csharp-spec.html}
\bibitem{src27} Referencia Normativa src-gamma-gof. \url{https://doc.rust-lang.org/reference/src-gamma-gof.html}
\bibitem{src28} Referencia Normativa src-kleppmann-ddia. \url{https://doc.rust-lang.org/reference/src-kleppmann-ddia.html}
\bibitem{src29} Referencia Normativa src-java-jls21. \url{https://doc.rust-lang.org/reference/src-java-jls21.html}
\bibitem{src30} Referencia Normativa src-kotlin-spec. \url{https://doc.rust-lang.org/reference/src-kotlin-spec.html}
\bibitem{src31} Referencia Normativa src-pierce-tapl. \url{https://doc.rust-lang.org/reference/src-pierce-tapl.html}
\end{thebibliography}
\end{document}
